\providecommand{\todo}[1]{\textcolor{red}{\textbf{[TODO: #1]}}}
\clearpage
\phantomsection
\section*{Chương 2. Chuẩn giao tiếp Avalon-ST, ngân sách chu kỳ khi xử lý tệp và kỹ thuật mã hóa FSM}
\addcontentsline{toc}{section}{Chương 2. Chuẩn giao tiếp Avalon-ST, ngân sách chu kỳ khi xử lý tệp và kỹ thuật mã hóa FSM}

\begin{learningobj}
\begin{itemize}[leftmargin=1.5em, itemsep=2pt]
    \item Thấu hiểu bản chất cơ chế bắt tay dòng chảy dữ liệu (streaming handshake) theo chuẩn giao tiếp công nghiệp Intel Avalon Streaming (Avalon-ST), phân biệt rõ ràng giữa hai tín hiệu kiểm soát \texttt{valid} (hợp lệ) và \texttt{ready} (sẵn sàng).
    \item Phân tích cặn kẽ cạm bẫy kỹ thuật của thông số \texttt{readyLatency} (độ trễ sẵn sàng 0 chu kỳ so với 1 chu kỳ), giải mã cơ chế gây thất thoát dữ liệu âm thanh khi kết nối các khối sở hữu trí tuệ (IP core) của Intel.
    \item Xác lập ngân sách chu kỳ xung nhịp chính xác từ nguyên lý đầu tiên cho hệ thống ước lượng cao độ $F_0$ khi dữ liệu đầu vào được cấp phát từ tệp WAV thay vì từ bộ giải mã âm thanh ngoại vi; phân định rõ sự khác biệt giữa thước đo thời gian thực vật lý và tốc độ tính toán phần cứng.
    \item Khảo sát toàn diện 7 phương án mã hóa máy trạng thái hữu hạn FSM (Finite State Machine) trên phần mềm tổng hợp Quartus Prime; đánh giá sự đánh đổi giữa tài nguyên thanh ghi flip-flop (FF) và cổng logic tổ hợp (LUT) trên kiến trúc FPGA Intel Cyclone V SoC.
    \item Chứng minh toán học chặt chẽ tính chất chuyển mạch đơn bit của mã Gray (\texttt{n \^{} (n >> 1)}) thông qua lập luận chuỗi bit 1 tận cùng (trailing-ones argument).
    \item Thiết kế và hiện thực hóa khối đệm trượt (Skid Buffer) hai tầng chốt bằng ngôn ngữ Verilog-2001 nhằm triệt tiêu đường phản áp tổ hợp ngược (combinational backpressure loop), bảo toàn 100\% thông lượng truyền thông.
    \item Xây dựng mô hình song sinh phần mềm (Python twin) và kịch bản kiểm thử chu kỳ (cycle-accurate testbench) với hơn 500 vector kích thích ngẫu nhiên, chứng minh hai phương án mã hóa FSM nhị phân và one-hot cho kết quả trùng khớp hoàn toàn.
\end{itemize}
\end{learningobj}

\subsection*{2.1. Chuẩn giao thức Avalon-ST: Cơ chế valid/ready và cạm bẫy readyLatency}
\addcontentsline{toc}{subsection}{2.1. Chuẩn giao thức Avalon-ST: Cơ chế valid/ready và cạm bẫy readyLatency}

\subsubsection*{Vấn đề thực tế: Dòng chảy dữ liệu không xung nhịp ép bước}
Khi thiết kế các hệ thống xử lý tín hiệu số (DSP -- Digital Signal Processing) trên vi mạch FPGA (Field-Programmable Gate Array), người lập trình thường bắt đầu với các mạch tuần tự đơn giản: mỗi nhịp xung nhịp đồng hồ, một mẫu dữ liệu mới đi vào, chạy qua các tầng thanh ghi và cho ra kết quả ở đầu bên kia. Mô hình đó hoạt động chính xác khi và chỉ khi nguồn phát (source) và nơi nhận (sink) luôn luôn sẵn sàng ở mọi nhịp xung nhịp.

Tuy nhiên, trong một hệ thống thực tế như thuật toán ước lượng cao độ tiếng nói $F_0$, dòng chảy dữ liệu có tính chất gián đoạn:
\begin{itemize}[leftmargin=1.5em, itemsep=2pt]
    \item Khối giải nén và phân phối mẫu âm thanh từ bộ nhớ hàng đợi FIFO (First-In First-Out) có thể phải tạm dừng một nhịp khi bộ đệm rỗng hoặc khi chuyển đổi giữa từ 32-bit sang mẫu 16-bit.
    \item Khối nhân tích lũy tự tương quan MAC (Multiply-Accumulate) cần tới $98.543$ chu kỳ xung nhịp để quét toàn bộ 189 độ trễ $\tau \in [26, 214]$ cho một khung 638 mẫu. Trong suốt thời gian tính toán này, khối MAC hoàn toàn không thể tiếp nhận thêm mẫu âm thanh mới của khung kế tiếp.
    \item Nếu khối phát đẩy mẫu lên bus mà không kiểm tra trạng thái bận của khối nhận, các mẫu âm thanh mới sẽ ghi đè lên mẫu cũ, làm vỡ cấu trúc khung và làm sai lệch quá trình tính toán.
\end{itemize}

Để giải quyết bài toán đồng bộ luồng dữ liệu giữa các khối có tốc độ xử lý chênh lệch, tập đoàn Intel chuẩn hóa giao diện truyền thông điểm-tới-điểm mang tên \textbf{Avalon Streaming Interface (Avalon-ST)}, được đặc tả chi tiết trong tài liệu kỹ thuật \emph{Intel Avalon Interface Specifications} (Mã tài liệu: MNL-AVALONIF).

\subsubsection*{Nguyên lý bắt tay valid/ready: Khi nào một mẫu được coi là chuyển giao?}
Trọng tâm của chuẩn Avalon-ST là cơ chế bắt tay hai chiều thông qua cặp tín hiệu điều khiển đơn bit:
\begin{enumerate}[leftmargin=1.5em, itemsep=2pt]
    \item \texttt{valid} (Phía phát đưa ra): Bằng 1 khi phía phát đã đặt một mẫu dữ liệu hợp lệ lên các đường dây dữ liệu \texttt{data}.
    \item \texttt{ready} (Phía nhận đưa ra): Bằng 1 khi phía nhận có đủ không gian bộ đệm và sẵn sàng nhận dữ liệu từ bus vào thanh ghi của mình.
\end{enumerate}

Nguyên lý cơ bản trong chuẩn giao tiếp là: \textbf{Một giao dịch truyền nhận dữ liệu (beat) chỉ được hoàn thành tại chu kỳ xung nhịp $n$ khi và chỉ khi cả hai tín hiệu \texttt{valid} và \texttt{ready} cùng bằng 1 tại chính chu kỳ đó}:
\begin{equation}
\text{Transfer}(n) = \texttt{valid}(n) \ \land \ \texttt{ready}(n)
\end{equation}

Nếu \texttt{valid} = 1 nhưng \texttt{ready} = 0, phía nhận đang áp dụng cơ chế \textbf{phản áp (backpressure)}. Khi đó, không có giao dịch nào diễn ra. Quy tắc bắt buộc đối với phía phát là: \textbf{phải giữ nguyên giá trị dữ liệu \texttt{data} và giữ \texttt{valid} = 1 ở chu kỳ kế tiếp}, phía phát không được thay đổi dữ liệu hoặc hạ \texttt{valid} xuống 0 cho đến khi phía nhận bật \texttt{ready} = 1 trở lại.

\subsubsection*{Cạm bẫy kỹ thuật: \texttt{readyLatency} = 0 so với \texttt{readyLatency} = 1}
Tưởng chừng nguyên lý trên đã đủ chặt chẽ, nhưng một cạm bẫy kỹ thuật nảy sinh khi ghép nối các khối phần cứng: \emph{Tín hiệu \texttt{ready} nói về khả năng tiếp nhận ở chính nhịp này hay ở nhịp sau?} Câu trả lời nằm ở tham số cấu hình mang tên \texttt{readyLatency} (độ trễ sẵn sàng).

Theo chuẩn giao tiếp Avalon-ST (\emph{Intel Avalon Interface Specifications}, Mục 3.4):
\begin{itemize}[leftmargin=1.5em, itemsep=3pt]
    \item \textbf{Cấu hình \texttt{readyLatency} = 0 (Bắt tay tức thời):} Quyết định nhận mẫu nằm ở \emph{chính chu kỳ hiện tại}. Nếu \texttt{ready} hạ xuống 0 ở chu kỳ $n$, mẫu đang nằm trên bus ở chu kỳ $n$ không được nhận, và phía phát phải tiếp tục giữ mẫu đó sang chu kỳ $n+1$. Ưu điểm của cấu hình này là phản hồi tức thì, không cần bộ đệm dự phòng. Nhược điểm là tạo ra một đường trễ tổ hợp (combinational path) nối thẳng từ khối nhận ngược về khối phát.
    \item \textbf{Cấu hình \texttt{readyLatency} = 1 (Tín hiệu sẵn sàng có trễ / cơ chế tín chỉ):} Tín hiệu \texttt{ready} ở chu kỳ $n$ là \emph{lời hứa cho chu kỳ $n+1$}. Khi khối nhận bật \texttt{ready} = 1 ở chu kỳ $n$, nó phát đi một tín chỉ: ``Ở chu kỳ $n+1$, tôi chắc chắn có chỗ để nhận một mẫu''. Ngược lại, nếu khối nhận hạ \texttt{ready} = 0 ở chu kỳ $n$, điều đó có nghĩa là ở chu kỳ $n+1$ khối phát không được gửi, nhưng \emph{mẫu xuất hiện ở chính chu kỳ $n$ vẫn được tiếp nhận đầy đủ}, bởi vì lời hứa tiếp nhận cho chu kỳ $n$ đã được khối nhận cam kết từ chu kỳ $n-1$.
\end{itemize}

\textbf{Cạm bẫy nối ghép IP:} Nhiều khối IP lõi xử lý DSP chuyên dụng của Intel (như bộ lọc FIR Compiler hay biến đổi Fourier nhanh FFT IP) mặc định thiết lập \texttt{readyLatency} = 1 trong cửa sổ cấu hình Parameter Editor nhằm chèn thêm một tầng thanh ghi vào đường tín hiệu \texttt{ready}, giúp tăng tần số hoạt động cực đại ($F_{\max}$). Trong khi đó, các kỹ sư viết mã Verilog bằng tay thường mặc định suy nghĩ theo kiểu \texttt{readyLatency} = 0.

Nếu khối phát giả định \texttt{readyLatency} = 1 trong khi khối nhận lại hoạt động với \texttt{readyLatency} = 0, sự cố mất đồng bộ dữ liệu sẽ xảy ra khi đường phản áp xuất hiện: một mẫu dữ liệu sẽ bị rơi mất khỏi dòng âm thanh mà không có bất kỳ cờ cảnh báo lỗi nào được bật lên.

\subsubsection*{Ví dụ tính tay và đối soát chu kỳ: Sự cố mất mẫu 12}
Để thấy rõ cơ chế phá hỏng dòng dữ liệu, hãy khảo sát một chuỗi truyền thông thực nghiệm gồm 4 mẫu âm thanh nguyên $x = [10, 11, 12, 13]$ đi qua kênh giao tiếp Avalon-ST từ chu kỳ 3 đến chu kỳ 8. Giả sử khối nhận gặp tắc nghẽn bộ đệm và bắt buộc phải hạ tín hiệu \texttt{ready} xuống 0 duy nhất tại chu kỳ 5 (\texttt{ready} lần lượt là $[1, 1, 0, 1, 1, 1]$).

Dữ liệu tính toán dưới đây được trích xuất trực tiếp từ kịch bản kiểm chứng mã nguồn\\ \texttt{verify\_v3/ch02\_avalon\_st\_trace.py}:

\begin{table}[htbp]
\centering
\footnotesize
\caption{Đối chiếu chu kỳ truyền nhận giữa trường hợp chuẩn và trường hợp lệch tham số \texttt{readyLatency}}
\label{tab:ch02_readylatency_trace}
\begin{tabular}{ccccccl}
\toprule
Chu kỳ & \texttt{valid} & \texttt{data} & \texttt{ready} & Bắt tay (RL=0) & Phía nhận lấy? & Hành vi thực tế của khối phát \\
\midrule
\multicolumn{7}{l}{\textbf{Kịch bản 1: Cả hai phía thống nhất cấu hình \texttt{readyLatency} = 0}} \\
3 & 1 & 10 & 1 & Có & Mẫu 10 & Khối phát đưa tiếp mẫu 11 lên bus. \\
4 & 1 & 11 & 1 & Có & Mẫu 11 & Khối phát đưa tiếp mẫu 12 lên bus. \\
5 & 1 & 12 & 0 & Không & Không & \texttt{ready}=0: Khối phát \textbf{giữ nguyên} mẫu 12. \\
6 & 1 & 12 & 1 & Có & Mẫu 12 & Khối phát đưa tiếp mẫu 13 lên bus. \\
7 & 1 & 13 & 1 & Có & Mẫu 13 & Đã gửi xong 4 mẫu, hạ \texttt{valid} về 0. \\
8 & 0 & -- & 1 & Không & Không & Bus nghỉ, nhận đủ 4 mẫu: $[10, 11, 12, 13]$. \\
\midrule
\multicolumn{7}{l}{\textbf{Kịch bản 2: Bất đồng cấu hình (Phát hiểu RL=1, Nhận hiểu RL=0)}} \\
3 & 1 & 10 & 1 & Có & Mẫu 10 & Khối phát gửi 10 do có tín chỉ từ CK 2. \\
4 & 1 & 11 & 1 & Có & Mẫu 11 & Khối phát gửi 11 do có tín chỉ từ CK 3. \\
5 & 1 & 12 & 0 & \textbf{Lệch} & \textbf{Từ chối} & Khối phát thấy CK4 \texttt{ready}=1 nên đẩy 12. \\
  &   &    &   &        &                & \emph{Nhưng nhận (RL=0) từ chối! Mẫu 12 bị rơi.} \\
6 & 0 & -- & 1 & Không  & Không          & Khối phát thấy CK5 \texttt{ready}=0 nên dừng phát. \\
7 & 1 & 13 & 1 & Có     & Mẫu 13         & Khối phát thấy CK6 \texttt{ready}=1 nên gửi 13. \\
8 & 0 & -- & 1 & Không  & Không          & Kết thúc: Khối nhận chỉ lấy $[10, 11, 13]$! \\
\bottomrule
\end{tabular}
\end{table}

\begin{figure}[htbp]
\centering
% chuong_v3/fig_ch02_readylatency.tex
% Comparison of Avalon-ST handshake with readyLatency = 0 versus readyLatency = 1
\begin{tikzpicture}[x=1.35cm, y=0.85cm, >=stealth, font=\small]

    % Style definitions
    \tikzstyle{siglabel} = [anchor=east, font=\footnotesize\ttfamily]
    \tikzstyle{gridline} = [dashed, color=gray!35, line width=0.5pt]
    \tikzstyle{sigline}  = [thick, color=blue!80!black]
    \tikzstyle{busline}  = [thick, color=teal!80!black]
    \tikzstyle{datalabel} = [font=\footnotesize\bfseries, color=teal!90!black]

    % Title 1: readyLatency = 0
    \node[anchor=west, font=\normalsize\bfseries, color=blue!70!black] at (-1.8, 4.8) 
        {(a) Giao tiếp Avalon-ST với \texttt{readyLatency} = 0 (Bắt tay cùng chu kỳ)};

    % Vertical cycle markers for Part (a)
    \foreach \c in {3,4,5,6,7,8} {
        \draw[gridline] (\c-3, 0.4) -- (\c-3, 4.3);
        \node[font=\footnotesize, color=gray!90!black] at (\c-3 + 0.5, 4.3) {CK \c};
    }
    \draw[gridline] (6, 0.4) -- (6, 4.3);

    % Signals for Part (a)
    % 1. clk
    \node[siglabel] at (-0.2, 3.6) {clk};
    \draw[thick, color=gray!70!black] 
        (0,3.3) \foreach \x in {0,1,2,3,4,5} { -- ++(0.5,0.6) -- ++(0.5,-0.6) };

    % 2. valid
    \node[siglabel] at (-0.2, 2.6) {valid};
    \draw[sigline] (0,2.9) -- (5,2.9) -- (5,2.3) -- (6,2.3);

    % 3. ready (drops at cycle 5)
    \node[siglabel] at (-0.2, 1.6) {ready};
    \draw[sigline] (0,1.9) -- (2,1.9) -- (2,1.3) -- (3,1.3) -- (3,1.9) -- (6,1.9);

    % 4. data bus
    \node[siglabel] at (-0.2, 0.7) {data};
    % Beat 10: [0, 1]
    \draw[busline] (0,0.5) -- (0,0.9) -- (1,0.9) -- (1,0.5) -- (0,0.5);
    \node[datalabel] at (0.5, 0.7) {10};
    % Beat 11: [1, 2]
    \draw[busline] (1,0.5) -- (1,0.9) -- (2,0.9) -- (2,0.5) -- (1,0.5);
    \node[datalabel] at (1.5, 0.7) {11};
    % Beat 12: [2, 4] (held across cycle 5!)
    \draw[busline] (2,0.5) -- (2,0.9) -- (4,0.9) -- (4,0.5) -- (2,0.5);
    \node[datalabel] at (3.0, 0.7) {12 (giữ do ready=0)};
    % Beat 13: [4, 5]
    \draw[busline] (4,0.5) -- (4,0.9) -- (5,0.9) -- (5,0.5) -- (4,0.5);
    \node[datalabel] at (4.5, 0.7) {13};
    % Idle: [5, 6]
    \draw[dashed, color=gray!60] (5,0.7) -- (6,0.7);

    % Transfer markers (checkmarks)
    \node[font=\scriptsize\bfseries, color=green!60!black] at (0.5, 0.15) {Nhận (10)};
    \node[font=\scriptsize\bfseries, color=green!60!black] at (1.5, 0.15) {Nhận (11)};
    \node[font=\scriptsize\bfseries, color=red!70!black]   at (2.5, 0.15) {Chờ (stall)};
    \node[font=\scriptsize\bfseries, color=green!60!black] at (3.5, 0.15) {Nhận (12)};
    \node[font=\scriptsize\bfseries, color=green!60!black] at (4.5, 0.15) {Nhận (13)};

    % Divider
    \draw[thick, color=gray!30] (-1.8, -0.4) -- (6.2, -0.4);

    % Title 2: readyLatency = 1
    \node[anchor=west, font=\normalsize\bfseries, color=purple!70!black] at (-1.8, -1.0) 
        {(b) Giao tiếp Avalon-ST với \texttt{readyLatency} = 1 (Tín hiệu ready báo trước 1 chu kỳ)};

    % Vertical cycle markers for Part (b)
    \foreach \c in {3,4,5,6,7,8} {
        \draw[gridline] (\c-3, -5.2) -- (\c-3, -1.4);
        \node[font=\footnotesize, color=gray!90!black] at (\c-3 + 0.5, -1.4) {CK \c};
    }
    \draw[gridline] (6, -5.2) -- (6, -1.4);

    % Signals for Part (b)
    % 1. clk
    \node[siglabel] at (-0.2, -2.1) {clk};
    \draw[thick, color=gray!70!black] 
        (0,-2.4) \foreach \x in {0,1,2,3,4,5} { -- ++(0.5,0.6) -- ++(0.5,-0.6) };

    % 2. ready (drops at cycle 5)
    \node[siglabel] at (-0.2, -3.1) {ready};
    \draw[thick, color=purple!80!black] (0,-2.8) -- (2,-2.8) -- (2,-3.4) -- (3,-3.4) -- (3,-2.8) -- (6,-2.8);

    % 3. valid
    \node[siglabel] at (-0.2, -4.0) {valid};
    % valid is 1 at CK3, CK4, CK5; drops at CK6 because ready was 0 at CK5!
    \draw[sigline] (0,-3.7) -- (3,-3.7) -- (3,-4.3) -- (4,-4.3) -- (4,-3.7) -- (5,-3.7) -- (5,-4.3) -- (6,-4.3);

    % 4. data bus
    \node[siglabel] at (-0.2, -4.8) {data};
    % Beat 10: [0, 1]
    \draw[busline] (0,-5.0) -- (0,-4.6) -- (1,-4.6) -- (1,-5.0) -- (0,-5.0);
    \node[datalabel] at (0.5, -4.8) {10};
    % Beat 11: [1, 2]
    \draw[busline] (1,-5.0) -- (1,-4.6) -- (2,-4.6) -- (2,-5.0) -- (1,-5.0);
    \node[datalabel] at (1.5, -4.8) {11};
    % Beat 12: [2, 3] (transferred at cycle 5 due to credit from cycle 4!)
    \draw[busline] (2,-5.0) -- (2,-4.6) -- (3,-4.6) -- (3,-5.0) -- (2,-5.0);
    \node[datalabel] at (2.5, -4.8) {12};
    % Stall at cycle 6: [3, 4]
    \draw[dashed, color=gray!60] (3,-4.8) -- (4,-4.8);
    % Beat 13: [4, 5] (transferred at cycle 7 due to credit from cycle 6!)
    \draw[busline] (4,-5.0) -- (4,-4.6) -- (5,-4.6) -- (5,-5.0) -- (4,-5.0);
    \node[datalabel] at (4.5, -4.8) {13};
    % Idle: [5, 6]
    \draw[dashed, color=gray!60] (5,-4.8) -- (6,-4.8);

    % Action notes with safe spacing
    \node[font=\tiny, color=purple!80!black] at (0.5, -5.45) {Tín chỉ CK2};
    \node[font=\tiny, color=purple!80!black] at (1.5, -5.45) {Tín chỉ CK3};
    \node[font=\tiny\bfseries, color=purple!80!black] at (2.5, -5.45) {Tín chỉ CK4};
    \node[font=\tiny\bfseries, color=red!70!black]   at (3.5, -5.45) {Dừng (CK5=0)};
    \node[font=\tiny, color=purple!80!black] at (4.5, -5.45) {Tín chỉ CK6};

\end{tikzpicture}

\caption{Dạng sóng thời gian so sánh cơ chế bắt tay Avalon-ST giữa cấu hình \texttt{readyLatency} = 0 và \texttt{readyLatency} = 1 khi xuất hiện xung phản áp ở chu kỳ 5 (tham chiếu \emph{Intel Avalon Interface Specifications}, Mục 3.4.2).}
\label{fig:ch02_readylatency}
\end{figure}

\textbf{Đọc hình.} Hình~\ref{fig:ch02_readylatency} trực quan hóa sự khác biệt bản chất giữa hai cấu hình:
\begin{itemize}[leftmargin=1.5em, itemsep=2pt]
    \item Ở nửa trên (a), với \texttt{readyLatency} = 0, tại chu kỳ 5 khi tín hiệu \texttt{ready} hạ xuống mức thấp, giao dịch dừng lại ngay lập tức. Mẫu 12 trên bus dữ liệu được giữ phẳng suốt hai chu kỳ xung nhịp (CK 5 và CK 6). Đến khi \texttt{ready} bật lên 1 ở CK 6, mẫu 12 mới được nhận và mẫu 13 nối tiếp ở CK 7. Không có mẫu nào bị mất.
    \item Ở nửa dưới (b), với \texttt{readyLatency} = 1, tín hiệu \texttt{ready} mang ý nghĩa tín chỉ trễ một nhịp. Vì ở chu kỳ 4 \texttt{ready} = 1, khối phát được quyền đẩy mẫu 12 ở chu kỳ 5 bất chấp việc tại chính chu kỳ 5 \texttt{ready} đang ở mức 0. Phía phát coi như mẫu 12 đã bàn giao thành công. Sang chu kỳ 6, khối phát mới nhận lệnh dừng do \texttt{ready} ở CK 5 bằng 0.
\end{itemize}

Hậu quả của sự bất đồng này là mẫu âm thanh số 12 bị loại khỏi dòng dữ liệu. Trong một khung phân tích 638 mẫu của thuật toán Praat, việc mất đúng một mẫu làm toàn bộ 637 mẫu phía sau bị dồn lên trước một vị trí thời gian. Các phép tính tự tương quan $R[\tau]$ sau đó vẫn thực thi bình thường theo công thức toán học, nhưng trên một cửa sổ tín hiệu bị lệch pha, dẫn đến sai lệch đỉnh parabol và xuất hiện lỗi thô (Gross Pitch Error) mà người lập trình khó tìm ra nguyên nhân nếu chỉ nhìn vào mã thuật toán.

\textit{Tóm lại: Tham số \texttt{readyLatency} quyết định thời điểm một mẫu dữ liệu được công nhận đã chuyển giao; hiểu lệch một nhịp giữa hai khối giao tiếp sẽ làm rơi hoặc nhân đôi mẫu âm thanh một cách im lặng mà không hề báo lỗi trên đường truyền.}

\subsection*{2.2. Ngân sách chu kỳ khi đầu vào là tệp âm thanh và thực tế thời gian thực}
\addcontentsline{toc}{subsection}{2.2. Ngân sách chu kỳ khi đầu vào là tệp âm thanh và thực tế thời gian thực}

\subsubsection*{Vấn đề thực tế: Nhận thức về đồng hồ lấy mẫu vật lý 16 kHz}
Khi tiếp cận bài toán xử lý tiếng nói trên FPGA, trực giác ban đầu của nhiều kỹ sư là tìm kiếm một bộ tạo xung nhịp 16 kHz trong thiết kế để điều nhịp cho toàn bộ hệ thống. Trực giác đó xuất phát từ thói quen làm việc với các bộ giải mã âm thanh tương tự (như chip WM8731 trên bo mạch DE10-Standard), nơi mà bộ chuyển đổi ADC (Analog-to-Digital Converter) liên tục xuất ra từng mẫu sau mỗi khoảng thời gian đều đặn $T_s = 1 / 16.000\text{ s} = 62{,}5\ \mu\text{s}$. Khi có một bộ ADC ngoại vi ép nhịp, nếu mạch FPGA xử lý chậm hơn $62{,}5\ \mu\text{s}$, bộ đệm sẽ tràn và dữ liệu âm thanh bị mất.

Tuy nhiên, trong kiến trúc hệ thống của tài liệu này (đã định hình ở Chương 1), \textbf{nguồn dữ liệu đầu vào là một tệp âm thanh WAV lưu trên thẻ nhớ SD, được vi xử lý HPS đọc và đẩy qua cầu Lightweight HPS-to-FPGA Bridge vào bộ đệm On-Chip FIFO}. 

Đặc điểm kỹ thuật của hệ thống là: \textbf{Không có đồng hồ 16 kHz nào ép bước trên phần cứng}.
\begin{itemize}[leftmargin=1.5em, itemsep=2pt]
    \item Vi xử lý HPS đẩy dữ liệu qua cầu giao tiếp Lightweight HPS-to-FPGA Bridge nhanh hơn rất nhiều so với tốc độ tiêu thụ 1 mẫu / 3.125 chu kỳ của chuỗi xử lý âm thanh (số chu kỳ ghi thực tế qua cầu truyền thông cần được đo đạc cụ thể trên phần cứng, như đã phân tích ở Chương 1). Nếu bộ đệm FIFO còn ô trống, HPS có thể ghi liên tục hàng nghìn mẫu vào phần cứng mà không bị nghẽn bởi tốc độ xung nhịp lấy mẫu vật lý.
    \item Lõi xử lý DSP trên FPGA chạy trên miền xung nhịp hệ thống $f_{\text{clk}} = 50\text{ MHz}$ (chu kỳ $T_{\text{clk}} = 20\text{ ns}$).
    \item Tốc độ tiêu thụ dữ liệu không bị khống chế bởi tần số lấy mẫu 16 kHz của âm thanh, mà bị quyết định trực tiếp bởi \textbf{thời gian thực thi của khối tính toán tự tương quan MAC}---khâu tiêu tốn nhiều thời gian tính toán nhất của toàn bộ hệ thống.
\end{itemize}

Vậy con số 16 kHz có ý nghĩa gì? Nó là thông số toán học quy định khoảng cách thời gian giữa các mẫu để tính ra cao độ thực theo đơn vị Hertz ($F_0 = f_s / \tau_{\text{interp}}$), và là \textbf{thước đo tham chiếu thời gian thực (real-time yardstick)} để đánh giá xem phần cứng chạy nhanh hơn thế giới thực bao nhiêu lần.

\subsubsection*{Dẫn xuất ngân sách chu kỳ từ nguyên lý đầu tiên}
Thiết lập ngân sách chu kỳ một cách tường minh từ các tham số chuẩn mực của phần mềm Praat 7.0.03 trên tệp âm thanh $f_s = 16.000\text{ Hz}$:
\begin{enumerate}[leftmargin=1.5em, itemsep=2pt]
    \item \textbf{Chu kỳ xung nhịp chuẩn:} Với tần số xung nhịp hệ thống $f_{\text{clk}} = 50\text{ MHz}$, mỗi chu kỳ xung nhịp kéo dài:
    \begin{equation}
    T_{\text{clk}} = \frac{1}{50.000.000\text{ Hz}} = 20\text{ ns}
    \end{equation}
    \item \textbf{Thước đo thời gian thực cho một mẫu:} Trong thế giới vật lý, một mẫu âm thanh 16 kHz xuất hiện sau mỗi khoảng thời gian $T_s = 62{,}5\ \mu\text{s}$. Số chu kỳ xung nhịp 50 MHz tương ứng với một khoảng thời gian lấy mẫu là:
    \begin{equation}
    C_{\text{sample}} = \frac{f_{\text{clk}}}{f_s} = \frac{50.000.000}{16.000} = 3.125\text{ chu kỳ/mẫu}
    \end{equation}
    \item \textbf{Bước nhảy khung phân tích (Frame Hop):} Thuật toán Praat với cấu hình mặc định (sàn cao độ $\text{pitchFloor} = 75\text{ Hz}$, bước thời gian tự động $\text{timeStep} = 0{,}75 / 75 = 10\text{ ms}$, tham chiếu mã nguồn \texttt{TODO-PRAAT-LINE}) phân tích các khung cách nhau đúng:
    \begin{equation}
    H = 10\text{ ms} \times 16.000\text{ mẫu/s} = 160\text{ mẫu}
    \end{equation}
    \item \textbf{Ngân sách chu kỳ thời gian thực cho một bước nhảy ($C_{\text{hop}}$):} Để đạt ngưỡng thời gian thực (xử lý kịp tốc độ nói của con người), toàn bộ chuỗi tính toán cho 160 mẫu mới phải hoàn thành trong khoảng thời gian không vượt quá 10 ms. Tại xung nhịp 50 MHz, ngân sách chu kỳ tối đa cho phép là:
    \begin{equation}
    C_{\text{hop}} = H \times C_{\text{sample}} = 160 \times 3.125 = 500.000\text{ chu kỳ/khung}
    \end{equation}
\end{enumerate}

\subsubsection*{Số phép tính MAC và chu kỳ thực thi của khối tự tương quan}
Mỗi khung âm thanh có độ dài cửa sổ $N = 638$ mẫu (đã được làm tròn từ $3 \times 16.000 / 75 = 640$ mẫu theo quy tắc đối xứng tâm của Boersma). Khối nhân tích lũy MAC cần tính toán:
\begin{itemize}[leftmargin=1.5em, itemsep=2pt]
    \item Năng lượng chuẩn $R[0]$: Cần tích lũy $638$ tích số $x_w[n] \times x_w[n]$ ($n = 0..637$). Số phép tính MAC là:
    \begin{equation}
    M_{R0} = 638\text{ phép tính}
    \end{equation}
    \item Dải quét độ trễ $\tau \in [26, 214]$ gồm 189 giá trị trễ liên tiếp: Với mỗi độ trễ $\tau$, số lượng tích số cần tính là $638 - \tau$. Khi $\tau = 26$, số phép nhân là $638 - 26 = 612$. Khi $\tau = 214$, số phép nhân là $638 - 214 = 424$. Tổng số phép tính MAC cho toàn bộ 189 độ trễ là tổng của một cấp số cộng:
    \begin{equation}
    M_{\text{lags}} = \sum_{\tau=26}^{214} (638 - \tau) = \frac{(612 + 424) \times 189}{2} = \frac{1.036 \times 189}{2} = 518 \times 189 = 97.902\text{ phép tính}
    \end{equation}
\end{itemize}

Tổng số phép tính nhân tích lũy trên một khung phân tích là:
\begin{equation}
M_{\text{total}} = M_{R0} + M_{\text{lags}} = 638 + 97.902 = 98.540\text{ phép tính}
\end{equation}

Do khối MAC được thiết kế dưới dạng đường ống dòng chảy liên tục (streaming pipeline, phát địa chỉ liền mạch không nhịp nghỉ như đã phân tích ở Chương 0), mỗi nhịp xung nhịp đưa vào một cặp toán hạng. Sau khi cặp toán hạng cuối cùng của $\tau = 214$ đi vào, bộ tích lũy trong khối DSP cần thêm đúng 3 chu kỳ trễ đường ống (pipeline latency) để đẩy kết quả cuối cùng ra thanh ghi xuất. Do đó, tổng số chu kỳ xung nhịp thực tế để tính xong toàn bộ ma trận tương quan của một khung là:
\begin{equation}
C_{\text{calc}} = M_{\text{total}} + 3 = 98.540 + 3 = 98.543\text{ chu kỳ/khung}
\end{equation}

Thời gian thực tế để FPGA xử lý xong một khung là:
\begin{equation}
T_{\text{calc}} = 98.543 \times 20\text{ ns} = 1.970.860\text{ ns} \approx 1{,}971\text{ ms}
\end{equation}

Chia số chu kỳ này cho 160 mẫu âm thanh mới nạp vào mỗi khung, ta có số chu kỳ tính toán trung bình trên một mẫu âm thanh:
\begin{equation}
C_{\text{per\_sample}} = \frac{98.543}{160} \approx 615{,}89\text{ chu kỳ/mẫu} \quad (\approx 616\text{ chu kỳ/mẫu})
\end{equation}

\subsubsection*{Định nghĩa và dẫn xuất Hệ số Thời gian thực (Real-Time Factor -- RTF)}
Hệ số thời gian thực $\text{RTF}$ là tỷ số không thứ nguyên giữa thời gian phần cứng thực thi phép tính ($T_{\text{calc}}$) và thời lượng âm thanh thực tế được xử lý ($T_{\text{audio}}$):
\begin{equation}
\text{RTF} = \frac{T_{\text{calc}}}{T_{\text{audio}}} = \frac{C_{\text{calc}} \times T_{\text{clk}}}{H \times T_s} = \frac{C_{\text{calc}}}{C_{\text{hop}}}
\end{equation}

Thay các con số đã dẫn xuất vào công thức:
\begin{equation}
\text{RTF}_{\text{MAC}} = \frac{98.543}{500.000} \approx 0{,}197086 \quad (\approx 0{,}197)
\end{equation}

Nếu tính thêm các khâu phụ trợ như trừ giá trị trung bình cục bộ, tra bảng Hanning, tìm đỉnh cực đại parabol và lọc trung vị (ước tính chiếm khoảng $1.157$ chu kỳ bổ sung, đưa tổng chu kỳ lõi lên mức $\approx 99.700$ chu kỳ, sẽ được chi tiết hóa ở Chương 12), hệ số thời gian thực toàn lõi là:
\begin{equation}
\text{RTF}_{\text{total}} = \frac{99.700}{500.000} \approx 0{,}1994 \quad (\approx 0{,}199)
\end{equation}

Con số $\text{RTF} \approx 0{,}197$ mang ý nghĩa vật lý rõ ràng: mạch FPGA xử lý tín hiệu nhanh hơn 5 lần so với thời gian thực ($1 / 0{,}197 \approx 5{,}07$). Một đoạn âm thanh đàm thoại dài 10 giây chỉ cần chưa đầy 2 giây để trích xuất toàn bộ đường nét cao độ $F_0$.

\subsubsection*{Cạm bẫy thường gặp: Nhầm lẫn giữa độ dài khung và bước nhảy thời gian}

\begin{tcolorbox}[colback=red!4!white, colframe=red!70!black, title=\textbf{Cạm bẫy thường gặp: Nhầm lẫn bước nhảy khung 10 ms và sự triệt tiêu sai số ngẫu nhiên}]
Khi thiết lập ngân sách chu kỳ cho hệ thống xử lý tiếng nói, người thiết kế rất dễ mắc phải hai cạm bẫy tính toán sau:
\begin{itemize}[leftmargin=1.2em, itemsep=2pt]
    \item \textbf{Cạm bẫy 1 (Nhầm lẫn bước nhảy khung $H$):} Nhầm lẫn giữa bước nhảy khung phân tích $H$ của thuật toán cao độ ($10\text{ ms}$, tương ứng 160 mẫu ở $16\text{ kHz}$) với bước nhảy $25\text{ ms}$ (400 mẫu) thường dùng trong các bài toán trích xuất đặc trưng MFCC. Tại xung nhịp 50 MHz, bước nhảy 25 ms cho ngân sách $1.250.000$ chu kỳ, lớn gấp 2,5 lần ngân sách chuẩn $500.000$ chu kỳ của Praat.
    \item \textbf{Cạm bẫy 2 (Sự triệt tiêu sai số ngẫu nhiên):} Nếu ai đó ước tính chu kỳ tính toán của một khung phân tích bằng một vòng lặp chưa tối ưu hóa với $\approx 246.000$ chu kỳ và chia cho ngân sách nhầm lẫn $1.250.000$ chu kỳ, tỷ số thu được là $246.000 / 1.250.000 \approx 0{,}197$. Con số này trùng hợp ngẫu nhiên với hệ số $\text{RTF} = 98.543 / 500.000 \approx 0{,}197$ chuẩn xác của chúng ta. Sự trùng hợp đó có thể làm người thiết kế ngộ nhận thiết kế đã đạt chuẩn, trong khi thực tế độ trễ phản hồi của hệ thống bị chậm đi 2,5 lần.
\end{itemize}
\end{tcolorbox}

Dưới đây là bảng đối chiếu tổng hợp được trích xuất trực tiếp từ kịch bản kiểm chứng số học\\ \texttt{verify\_v3/ch02\_cycle\_budget.py}:

\begin{table}[htbp]
\centering
\footnotesize
\setlength{\tabcolsep}{3.5pt}
\caption{Bảng đối chiếu thông số ngân sách chu kỳ và phân tích sai lệch kỹ thuật}
\label{tab:ch02_cycle_budget}
\begin{tabularx}{\textwidth}{p{4.2cm}ccX}
\toprule
Thông số hệ thống & Chuẩn Praat (Hop 10 ms) & Nhầm lẫn (Hop 25 ms) & Bản chất kỹ thuật \\
\midrule
Tần số lấy mẫu âm thanh ($f_s$) & 16.000 Hz & 16.000 Hz & Tiêu chuẩn âm thanh tiếng nói \\
Tần số xung nhịp FPGA ($f_{\text{clk}}$) & 50.000.000 Hz & 50.000.000 Hz & Xung nhịp đồng bộ hệ thống \\
Thước đo chu kỳ trên mẫu & 3.125 chu kỳ/mẫu & 3.125 chu kỳ/mẫu & $f_{\text{clk}} / f_s = 50\text{ MHz} / 16\text{ kHz}$ \\
Bước nhảy khung ($H$) & \textbf{160 mẫu (10 ms)} & \textbf{400 mẫu (25 ms)} & Praat dùng bước nhảy 10 ms; nhầm sang 25 ms \\
Ngân sách thời gian thực ($C_{\text{hop}}$) & \textbf{500.000 chu kỳ} & \textbf{1.250.000 chu kỳ} & $H \times 3.125$ chu kỳ \\
Độ dài cửa sổ phân tích ($N$) & 638 mẫu & 640 mẫu & Praat làm tròn tâm đối xứng \\
Số phép tính MAC ($M_{\text{total}}$) & 98.540 & Chưa tối ưu & $638 + \sum_{26}^{214} (638 - \tau)$ \\
Chu kỳ thực tế của lõi ($C_{\text{calc}}$) & \textbf{98.543 chu kỳ} & \textbf{246.000 chu kỳ} & 98.540 nhịp phát + 3 nhịp trễ \\
Chu kỳ tính toán trên 1 mẫu mới & \textbf{615,89 chu kỳ/mẫu} & 615,00 chu kỳ/mẫu & $C_{\text{calc}} / H$ \\
Thời gian tính toán một khung & \textbf{1,971 ms} & 4,920 ms & $C_{\text{calc}} \times 20\text{ ns}$ \\
Hệ số thời gian thực ($\text{RTF}_{\text{MAC}}$) & \textbf{0,197} & \textbf{0,197} & Trùng hợp ngẫu nhiên do sai cả tử và mẫu \\
\bottomrule
\end{tabularx}
\end{table}

\textbf{Đọc bảng.} Bảng~\ref{tab:ch02_cycle_budget} khẳng định nguyên tắc thiết kế quan trọng: thiết kế chuẩn xác hoàn tất một khung phân tích chỉ trong $1{,}97\text{ ms}$ trên bước nhảy $10\text{ ms}$ (đúng chuẩn tương quan của Boersma), giúp hệ thống có độ trễ phản hồi thấp và giải phóng hơn 80\% băng thông xử lý của FPGA cho các tác vụ tiếp theo.

\textit{Tóm lại: Khi xử lý tệp âm thanh, không có đồng hồ 16 kHz nào chi phối bước chạy; lõi FPGA tính toán ở tần số xung nhịp 50 MHz với 98.543 chu kỳ cho mỗi bước nhảy 160 mẫu (616 chu kỳ/mẫu), hoàn tất một khung 10 ms chỉ trong 1,97 ms, đạt hệ số thời gian thực RTF = 0,197 so với ngân sách 500.000 chu kỳ.}

\subsection*{2.3. Kỹ thuật mã hóa FSM: Đánh đổi tài nguyên và độ trễ trên kiến trúc Cyclone V}
\addcontentsline{toc}{subsection}{2.3. Kỹ thuật mã hóa FSM: Đánh đổi tài nguyên và độ trễ trên kiến trúc Cyclone V}

\subsubsection*{Vấn đề thực tế: Từ câu lệnh switch-case trừu tượng đến cổng logic vật lý}
Trong ngôn ngữ mô tả phần cứng Verilog, máy trạng thái hữu hạn FSM (Finite State Machine) thường được lập trình viên biểu diễn dưới dạng câu lệnh rẽ nhánh trừu tượng: \texttt{case (state\_reg) ... endcase}. Nhìn trên màn hình soạn thảo mã nguồn, việc chuyển trạng thái trông giống như một phép gán biến tức thời trong ngôn ngữ C.

Nhưng lớp silicon bên dưới không có khái niệm ``trạng thái trừu tượng''. Mọi trạng thái bắt buộc phải được số hóa thành một chuỗi các bit nhị phân lưu trữ trong các phần tử nhớ tuần tự---các thanh ghi flip-flop (FF). Phần mềm tổng hợp logic Quartus Prime khi phân tích mã nguồn sẽ phải đối mặt với câu hỏi cốt lõi: \textbf{Mỗi trạng thái nên được gán bằng tổ hợp bit nào để mạch chạy nhanh nhất ($F_{\max}$ cao nhất) mà tốn ít diện tích nhất?}

Cách thức gán bit cho trạng thái (State Encoding) chi phối trực tiếp hai yếu tố của mạch tích hợp:
\begin{enumerate}[leftmargin=1.5em, itemsep=2pt]
    \item \textbf{Số lượng thanh ghi flip-flop tiêu tốn:} Ảnh hưởng tới dung lượng tài nguyên tuần tự của chip.
    \item \textbf{Độ sâu và độ phức tạp của mạng logic tổ hợp (Combinational Logic Depth):} Logic chuyển trạng thái tiếp theo (\texttt{next\_state}) và logic sinh tín hiệu ngõ ra (\texttt{outputs}) đòi hỏi mạng cổng logic LUT (Look-Up Table) có số đầu vào (fan-in) rộng hay hẹp. Mạng logic càng sâu, thời gian truyền trễ (propagation delay) càng lớn, kéo tụt tần số hoạt động tối đa $F_{\max}$ của hệ thống.
\end{enumerate}

\subsubsection*{Khảo sát 7 phương án mã hóa trong thiết lập State Machine Processing của Quartus}
Trong cửa sổ cấu hình dự án của Quartus Prime (\emph{Assignments $\to$ Settings $\to$ Compiler Settings $\to$ Advanced Settings (Synthesis) $\to$ State Machine Processing}, tham chiếu tài liệu \emph{Intel Quartus Prime Standard Edition User Guide: Design Compilation}), công cụ cung cấp đúng 7 phương án xử lý máy trạng thái:

\begin{enumerate}[leftmargin=1.5em, itemsep=3pt]
    \item \textbf{Auto (Tự động - Mặc định):} Quartus tự động phân tích cấu trúc FSM dựa trên thuật toán tối ưu hóa nội bộ. Với hầu hết các máy trạng thái trong FPGA có số trạng thái nhỏ hơn hoặc bằng 32, Quartus sẽ tự động chuyển đổi sang mã hóa One-Hot nhằm tối ưu hóa tần số $F_{\max}$.
    \item \textbf{One-Hot (Một bit tích cực):} Nếu FSM có $S$ trạng thái, bộ mã hóa sử dụng đúng $S$ thanh ghi flip-flop. Tại một thời điểm bất kỳ, chỉ có duy nhất 1 thanh ghi mang giá trị 1, tất cả các thanh ghi còn lại mang giá trị 0.
    \item \textbf{Minimal Bits (Nhị phân tối thiểu):} Số thanh ghi flip-flop sử dụng là giá trị nhỏ nhất đủ để biểu diễn $S$ trạng thái: $K = \lceil \log_2 S \rceil$. Các trạng thái được gán mã nhị phân nhằm tối thiểu hóa số bit và tối ưu hóa diện tích logic.
    \item \textbf{Sequential (Nhị phân tuần tự theo thứ tự khai báo):} Cũng sử dụng $K = \lceil \log_2 S \rceil$ thanh ghi flip-flop, nhưng Quartus gán các giá trị nhị phân tuần tự tăng dần ($0, 1, 2, \dots$) đúng theo thứ tự khai báo của các trạng thái trong mã nguồn HDL, không sắp xếp lại theo đồ thị chuyển trạng thái.
    \item \textbf{Gray (Mã Gray chuyển mạch đơn bit):} Sử dụng $K = \lceil \log_2 S \rceil$ thanh ghi, trong đó các trạng thái liền kề nhau trong chu trình hoạt động chỉ khác nhau đúng 1 bit duy nhất.
    \item \textbf{Johnson (Mã vòng lặp đảo Moebius):} Sử dụng $K = \lceil S / 2 \rceil$ thanh ghi dưới dạng một thanh ghi dịch vòng tròn với đầu ra đảo quay về đầu vào. Hai trạng thái liên tiếp cũng chỉ lệch nhau đúng 1 bit.
    \item \textbf{User-Encoded (Mã hóa do người dùng định nghĩa):} Quartus giữ nguyên các giá trị gán bit mà người thiết kế đã khai báo trong mã nguồn Verilog (\texttt{parameter} hoặc \texttt{localparam}), không tự ý trích xuất hoặc thay đổi.
\end{enumerate}

\subsubsection*{Phân tích so sánh trên máy trạng thái 5 trạng thái điển hình}
Để làm sáng tỏ bản chất của từng phương án, hãy xét một máy trạng thái điều khiển khung điển hình gồm $S = 5$ trạng thái xuất hiện trong bộ xử lý âm thanh:
\begin{itemize}[leftmargin=1.5em, itemsep=1pt]
    \item \texttt{S\_IDLE} (Trạng thái nghỉ, chờ tín hiệu \texttt{start})
    \item \texttt{S\_LOAD} (Nạp mẫu âm thanh từ FIFO vào RAM khung)
    \item \texttt{S\_CALC} (Kích hoạt lõi tính toán tự tương quan MAC)
    \item \texttt{S\_DRAIN} (Đẩy kết quả cao độ ra bộ đệm xuất)
    \item \texttt{S\_DONE} (Phát xung báo hoàn tất và quay về \texttt{S\_IDLE})
\end{itemize}

Dưới đây là bảng phân bổ tổ hợp bit và đánh giá chi phí phần cứng của 4 phương án mã hóa tiêu biểu:

\begin{table}[htbp]
\centering
\footnotesize
\setlength{\tabcolsep}{4.5pt}
\caption{So sánh các phương án mã hóa trạng thái cho FSM 5 trạng thái}
\label{tab:ch02_fsm_encodings}
\begin{tabular}{lcccc}
\toprule
Trạng thái FSM & Minimal Bits (Binary) & One-Hot & Mã Gray & Mã Johnson \\
\midrule
\texttt{S\_IDLE}  & \texttt{000} & \texttt{00001} & \texttt{000} & \texttt{000} \\
\texttt{S\_LOAD}  & \texttt{001} & \texttt{00010} & \texttt{001} & \texttt{001} \\
\texttt{S\_CALC}  & \texttt{010} & \texttt{00100} & \texttt{011} & \texttt{011} \\
\texttt{S\_DRAIN} & \texttt{011} & \texttt{01000} & \texttt{010} & \texttt{111} \\
\texttt{S\_DONE}  & \texttt{100} & \texttt{10000} & \texttt{110} & \texttt{110} \\
\midrule
Số flip-flop ($K$) & \textbf{3 FF} & \textbf{5 FF} & \textbf{3 FF} & \textbf{3 FF} \\
Số trạng thái thừa & $2^3 - 5 = 3$ & $2^5 - 5 = 27$ & $2^3 - 5 = 3$ & $2^3 - 5 = 3$ \\
Fan-in giải mã ngõ ra & 3 biến (cần LUT3) & \textbf{1 biến (Dây nối)} & 3 biến (cần LUT3) & 2 biến (cần LUT2) \\
Độ trễ giải mã ngõ ra & Trễ qua cổng LUT & \textbf{0 ns (Không tốn cổng)} & Trễ qua cổng LUT & Trễ qua cổng LUT \\
Chuyển mạch đơn bit? & Không (001 $\to$ 010: 2 bit) & Không (2 bit đổi) & \textbf{Có (luôn 1 bit)} & \textbf{Có (luôn 1 bit)} \\
\bottomrule
\end{tabular}
\end{table}

\textbf{Bản chất đánh đổi:}
\begin{itemize}[leftmargin=1.5em, itemsep=2pt]
    \item \textbf{Minimal Bits (Binary):} Tiết kiệm thanh ghi nhất (chỉ tốn 3 FF cho 5 trạng thái). Nhưng để biết máy trạng thái có đang ở \texttt{S\_CALC} hay không, mạch phải giải mã biểu thức logic tổ hợp 3 biến: $\texttt{calc\_en} = \overline{s_2} \cdot s_1 \cdot \overline{s_0}$. Tín hiệu này buộc phải đi qua một bảng tra LUT 3 đầu vào, kỳ vọng sẽ tốn thêm logic ALM và có thể thêm độ trễ vào đường tới hạn, cần được đo đạc qua báo cáo định thời của Quartus Fitter.
    \item \textbf{One-Hot:} Tốn nhiều thanh ghi nhất (5 FF cho 5 trạng thái). Nhưng đổi lại, để kích hoạt tín hiệu tính toán, ta chỉ cần nối thẳng đường dây từ flip-flop thứ 2: $\texttt{calc\_en} = s[2]$. Cổng giải mã ngõ ra không cần bảng tra LUT. Đường chuyển trạng thái kế tiếp cũng chỉ là các cổng AND/OR 2 đầu vào đơn giản ($s_{\text{next}}[2] = s[1] \cdot \texttt{load\_done} + s[2] \cdot \overline{\texttt{calc\_done}}$). Do đó, One-Hot giảm đáng kể độ trễ tổ hợp và cho tần số $F_{\max}$ thuận lợi.
\end{itemize}

\subsubsection*{Mã nguồn HDL: Hiện thực hóa FSM nhị phân và One-Hot}
Để phục vụ việc kiểm chứng và so sánh, hai module FSM tương đương chức năng được viết bằng Verilog-2001 chuẩn:

\begin{lstlisting}[language=Verilog, firstnumber=40, caption={Mã nguồn FSM 5 trạng thái mã hóa nhị phân Minimal Bits (rtl/v3/ch02\_fsm\_binary.v)}]
    localparam [2:0] S_IDLE  = 3'b000;
    localparam [2:0] S_LOAD  = 3'b001;
    localparam [2:0] S_CALC  = 3'b010;
    localparam [2:0] S_DRAIN = 3'b011;
    localparam [2:0] S_DONE  = 3'b100;

    reg [2:0] state_reg, state_next;

    // 1. State register
    always @(posedge clk or negedge rst_n) begin
        if (!rst_n) begin
            state_reg <= S_IDLE;
        end else begin
            state_reg <= state_next;
        end
    end

    // 2. Next-state combinational logic
    always @(*) begin
        state_next = state_reg;
        case (state_reg)
            S_IDLE: begin
                if (start) state_next = S_LOAD;
            end
            S_LOAD: begin
                if (load_done) state_next = S_CALC;
            end
            S_CALC: begin
                if (calc_done) state_next = S_DRAIN;
            end
            S_DRAIN: begin
                if (drain_done) state_next = S_DONE;
            end
            S_DONE: begin
                state_next = S_IDLE;
            end
            default: begin
                state_next = S_IDLE;
            end
        endcase
    end

    // 3. Moore outputs decoded directly from state_reg
    assign busy       = (state_reg != S_IDLE);
    assign load_en    = (state_reg == S_LOAD);
    assign calc_en    = (state_reg == S_CALC);
    assign drain_en   = (state_reg == S_DRAIN);
    assign done_pulse = (state_reg == S_DONE);
    assign state_num  = state_reg;
\end{lstlisting}

\textbf{Giải thích mã nguồn \texttt{ch02\_fsm\_binary.v} dòng-từng-dòng:}
\begin{itemize}[leftmargin=1.5em, itemsep=2pt]
    \item Dòng 40--44: Khai báo các hằng số trạng thái \texttt{localparam [2:0]} theo mã nhị phân từ \texttt{S\_IDLE} (0) đến \texttt{S\_DONE} (4).
    \item Dòng 46: Khai báo hai thanh ghi trạng thái 3-bit \texttt{state\_reg} (lưu trữ) và \texttt{state\_next} (dây tổ hợp).
    \item Dòng 49--55: Khối tuần tự \texttt{always @(posedge clk or negedge rst\_n)} cập nhật trạng thái tại cạnh lên xung nhịp, reset bất đồng bộ tích cực mức thấp về \texttt{S\_IDLE}.
    \item Dòng 58--80: Khối tổ hợp giải mã trạng thái tiếp theo qua câu lệnh \texttt{case (state\_reg)} từ các tín hiệu \texttt{start}, \texttt{load\_done}, \texttt{calc\_done}, \texttt{drain\_done}.
    \item Dòng 83--88: Các phép gán Moore: cờ \texttt{busy} tích cực khi khác \texttt{S\_IDLE}; các cờ điều khiển \texttt{load\_en}, \texttt{calc\_en}, \texttt{drain\_en}, \texttt{done\_pulse} giải mã qua phép so sánh bằng; và xuất chỉ số \texttt{state\_num = state\_reg}.
\end{itemize}

Tiếp theo là mã nguồn hiện thực hóa theo phương án One-Hot tường minh:

\begin{lstlisting}[language=Verilog, firstnumber=41, caption={Mã nguồn FSM 5 trạng thái mã hóa One-Hot tường minh (rtl/v3/ch02\_fsm\_onehot.v)}]
    localparam [4:0] S_IDLE  = 5'b00001;
    localparam [4:0] S_LOAD  = 5'b00010;
    localparam [4:0] S_CALC  = 5'b00100;
    localparam [4:0] S_DRAIN = 5'b01000;
    localparam [4:0] S_DONE  = 5'b10000;

    reg [4:0] state_reg, state_next;

    // 1. State register (5 flip-flops)
    always @(posedge clk or negedge rst_n) begin
        if (!rst_n) begin
            state_reg <= S_IDLE;
        end else begin
            state_reg <= state_next;
        end
    end

    // 2. One-hot next-state logic
    always @(*) begin
        state_next = 5'b00000;
        state_next[0] = (state_reg[0] & ~start) | state_reg[4];
        state_next[1] = (state_reg[0] & start)  | (state_reg[1] & ~load_done);
        state_next[2] = (state_reg[1] & load_done) | (state_reg[2] & ~calc_done);
        state_next[3] = (state_reg[2] & calc_done) | (state_reg[3] & ~drain_done);
        state_next[4] = (state_reg[3] & drain_done);
    end

    // 3. Moore outputs: single-bit wires directly from state flip-flops
    assign busy       = ~state_reg[0];
    assign load_en    = state_reg[1];
    assign calc_en    = state_reg[2];
    assign drain_en   = state_reg[3];
    assign done_pulse = state_reg[4];

    // Priority decode to numeric state number (0..4) for comparison
    always @(*) begin
        case (1'b1)
            state_reg[0]: state_num = 3'd0;
            state_reg[1]: state_num = 3'd1;
            state_reg[2]: state_num = 3'd2;
            state_reg[3]: state_num = 3'd3;
            state_reg[4]: state_num = 3'd4;
            default:      state_num = 3'd0;
        endcase
    end
\end{lstlisting}

\textbf{Giải thích mã nguồn \texttt{ch02\_fsm\_onehot.v} dòng-từng-dòng:}
\begin{itemize}[leftmargin=1.5em, itemsep=2pt]
    \item Dòng 41--45: Khai báo 5 hằng số trạng thái 5-bit, mỗi trạng thái có duy nhất 1 bit bật lên 1 (bit 0 cho \texttt{S\_IDLE}, bit 4 cho \texttt{S\_DONE}).
    \item Dòng 47: Khai báo thanh ghi trạng thái 5-bit \texttt{state\_reg} gồm 5 flip-flop vật lý.
    \item Dòng 50--56: Khối tuần tự chốt 5 flip-flop tại cạnh lên xung nhịp, reset về \texttt{S\_IDLE} (mã \texttt{5'b00001}).
    \item Dòng 59--66: Khối tổ hợp xác định từng bit của \texttt{state\_next} thông qua các biểu thức AND/OR 2 đầu vào, ví dụ: bit 2 bật lên khi vừa xong nạp mẫu (\texttt{state\_reg[1] \& load\_done}) hoặc đang duy trì tính toán (\texttt{state\_reg[2] \& \~{}calc\_done}).
    \item Dòng 69--73: Các tín hiệu ngõ ra Moore là các đường dây nối trực tiếp từ flip-flop của \texttt{state\_reg} (\texttt{load\_en} nối từ \texttt{state\_reg[1]}, \texttt{calc\_en} từ \texttt{state\_reg[2]}), không cần cổng logic giải mã.
    \item Dòng 76--85: Khối giải mã ưu tiên ánh xạ từ mã One-Hot sang chỉ số nguyên 3-bit \texttt{state\_num} (0 đến 4) để phục vụ kiểm chứng đối đầu với module nhị phân.
\end{itemize}

\textbf{Tương tác giữa trình tổng hợp Quartus và mã hóa One-Hot viết tay:}
Mặc định trong Quartus Prime, thiết lập \emph{State Machine Processing} được đặt là \textbf{Auto}. Bộ trích xuất máy trạng thái (State Machine Extractor) sẽ nhận diện khối tuần tự FSM và tự động mã hóa lại (re-encode) đồ thị trạng thái theo thuật toán tối ưu của công cụ.
Khi người thiết kế đã chủ động mã hóa One-Hot bằng tay trong mã nguồn Verilog:
\begin{itemize}[leftmargin=1.5em, itemsep=2pt]
    \item Nếu để \textbf{Auto}, Quartus vẫn có thể nhận diện FSM và tự động gán lại mã hóa theo thuật toán của mình.
    \item Để buộc Quartus giữ nguyên chính xác tổ hợp bit và cấu trúc phần cứng mà người thiết kế đã định nghĩa, cần thiết lập \emph{State Machine Processing} thành \textbf{User-Encoded}, hoặc thêm thuộc tính tổng hợp trực tiếp vào khai báo thanh ghi trong Verilog:
    \begin{equation}
    \texttt{(* syn\_encoding = "user" *) reg [4:0] state\_reg;}
    \end{equation}
    Thuộc tính \texttt{syn\_encoding = "user"} chỉ thị rõ ràng cho trình tổng hợp Quartus không được tự ý tối ưu hóa lại bảng mã hóa, giúp bảo toàn toàn bộ cấu trúc dây nối trực tiếp của thiết kế One-Hot.
\end{itemize}

\subsubsection*{Bối cảnh vi kiến trúc Cyclone V: Hiệu quả tài nguyên của mã hóa One-Hot}
Một quan niệm thường gặp khi học thiết kế số là hạn chế mã hóa One-Hot vì sợ tốn thanh ghi flip-flop, ưu tiên Minimal Bits. Tuy nhiên, quan niệm này không phản ánh đúng cấu trúc phần cứng hiện đại trên FPGA Intel Cyclone V SoC.

Trên bo mạch DE10-Standard, vi mạch FPGA được trang bị là Cyclone V SoC 5CSXFC6D6F31C6 (tham chiếu tài liệu chính thức \emph{Intel Cyclone V Device Overview}, mã tài liệu CV-51001, Bảng 1 ``Cyclone V SoC Family Plan''). Theo tài liệu kỹ thuật này:
\begin{itemize}[leftmargin=1.5em, itemsep=2pt]
    \item Vi mạch sở hữu đúng \textbf{41.910 khối logic thích nghi ALM (Adaptive Logic Module)}.
    \item Mỗi khối ALM của Cyclone V chứa các bảng tra LUT phân chia được (fracturable 8-input LUT) và tích hợp sẵn \textbf{bốn thanh ghi flip-flop chuyên dụng} (4 dedicated registers).
    \item Tổng số lượng thanh ghi flip-flop khả dụng trên toàn chip là:
    \begin{equation}
    N_{\text{registers}} = 41.910 \times 4 = \textbf{166.036 thanh ghi flip-flop}
    \end{equation}
\end{itemize}

Toàn bộ hệ thống ước lượng cao độ $F_0$ từ đầu đến cuối dự kiến sử dụng chưa tới $2.500$ thanh ghi (chiếm khoảng $1{,}5\%$ tài nguyên thanh ghi của chip). 

Trong một con chip có sẵn $166.036$ thanh ghi, việc sử dụng thêm 2 thanh ghi FF để chuyển FSM từ 3 bit (Binary) sang 5 bit (One-Hot) không gây áp lực về dung lượng tài nguyên tuần tự. Ngược lại, việc giảm được số tầng logic LUT tổ hợp sẽ giúp thiết kế dễ dàng đạt thời gian định thời (timing closure) ở 50 MHz mà không gây áp lực lên khâu định tuyến (routing) của phần mềm Quartus Fitter.

\subsubsection*{Chứng minh toán học: Bản chất chuyển mạch đơn bit của mã Gray}
Mã Gray (Gray Code) do Frank Gray công bố năm 1953, giữ vai trò quan trọng trong các bộ đếm con trỏ FIFO bất đồng bộ giữa hai miền xung nhịp khác nhau. Thuật toán chuyển đổi từ số nhị phân $n$ sang mã Gray $G(n)$ được hiện thực hóa bằng biểu thức:
\begin{equation}
G(n) = n \ \oplus \ (n \gg 1)
\end{equation}

Câu hỏi toán học mang tính nguyên lý đầu tiên là: \textbf{Vì sao công thức $n \oplus (n \gg 1)$ luôn đảm bảo rằng khi $n$ tăng lên $n+1$, giá trị mã Gray $G(n)$ và $G(n+1)$ chỉ thay đổi đúng duy nhất 1 bit?}

Dưới đây là chứng minh toán học thông qua \textbf{Lập luận chuỗi bit 1 tận cùng (Trailing-Ones Argument)}:

\textbf{Bước 1: Phân tích bản chất phép cộng 1 trong hệ nhị phân.}
Mọi số nguyên không âm $n$ khi biểu diễn dưới dạng nhị phân đều có thể viết dưới dạng tổng quát với $k$ bit 1 liên tiếp ở phần tận cùng bên phải ($k \ge 0$):
\begin{equation}
n = (\dots b_{k+1} \, 0 \, \underbrace{1 \, 1 \, \dots \, 1}_{k \text{ bit 1}})_2
\end{equation}
Khi ta cộng thêm 1 vào số $n$, quy tắc nhớ (carry propagation) của phép cộng nhị phân sẽ lật toàn bộ $k$ bit 1 này thành số 0, và phát sinh một bit nhớ làm lật bit thứ $k$ (vốn đang là 0) thành 1. Toàn bộ các bit từ vị trí $k+1$ trở về trước hoàn toàn không bị ảnh hưởng:
\begin{equation}
n + 1 = (\dots b_{k+1} \, 1 \, \underbrace{0 \, 0 \, \dots \, 0}_{k \text{ bit 0}})_2
\end{equation}

Gọi $\Delta_j = b_j(n+1) \oplus b_j(n)$ là chỉ số báo hiệu bit thứ $j$ của số nhị phân có bị đổi trạng thái hay không:
\begin{itemize}[leftmargin=1.5em, itemsep=2pt]
    \item Với các vị trí $0 \le j < k$: Bit cũ là 1, bit mới là 0 $\implies \Delta_j = 1 \oplus 0 = 1$.
    \item Với vị trí $j = k$: Bit cũ là 0, bit mới là 1 $\implies \Delta_k = 0 \oplus 1 = 1$.
    \item Với các vị trí $j > k$: Bit không đổi $\implies \Delta_j = 0$.
\end{itemize}

\textbf{Bước 2: Biểu diễn bit của mã Gray.}
Theo định nghĩa $G(n) = n \oplus (n \gg 1)$, bit thứ $j$ của mã Gray $g_j(n)$ được tính bằng phép XOR giữa hai bit liền kề của số nhị phân:
\begin{equation}
g_j(n) = b_j(n) \ \oplus \ b_{j+1}(n)
\end{equation}

Xét sự thay đổi của bit thứ $j$ giữa $G(n+1)$ và $G(n)$, ký hiệu là $\delta_j = g_j(n+1) \oplus g_j(n)$:
\begin{align}
\delta_j &= \big[b_j(n+1) \oplus b_{j+1}(n+1)\big] \ \oplus \ \big[b_j(n) \oplus b_{j+1}(n)\big] \nonumber \\
         &= \big[b_j(n+1) \oplus b_j(n)\big] \ \oplus \ \big[b_{j+1}(n+1) \oplus b_{j+1}(n)\big] \nonumber \\
         &= \Delta_j \ \oplus \ \Delta_{j+1}
\end{align}

\textbf{Bước 3: Đánh giá giá trị $\delta_j$ theo ba khoảng chỉ số $j$.}
\begin{enumerate}[leftmargin=1.5em, itemsep=2pt]
    \item \textbf{Khoảng $0 \le j < k$:} Tại khoảng này, cả $\Delta_j$ và $\Delta_{j+1}$ đều bằng 1 (vì $j+1 \le k$). Do đó:
    \begin{equation}
    \delta_j = 1 \ \oplus \ 1 = 0 \quad (\text{Bit thứ } j \text{ của mã Gray \textbf{không đổi}!})
    \end{equation}
    \item \textbf{Tại chính vị trí $j = k$:} Ta có $\Delta_k = 1$ và $\Delta_{k+1} = 0$. Do đó:
    \begin{equation}
    \delta_k = 1 \ \oplus \ 0 = 1 \quad (\text{Bit thứ } k \text{ của mã Gray \textbf{thay đổi}!})
    \end{equation}
    \item \textbf{Khoảng $j > k$:} Cả $\Delta_j$ và $\Delta_{j+1}$ đều bằng 0. Do đó:
    \begin{equation}
    \delta_j = 0 \ \oplus \ 0 = 0 \quad (\text{Bit thứ } j \text{ của mã Gray \textbf{không đổi}!})
    \end{equation}
\end{enumerate}

\textbf{Kết luận toán học:} Khi chuyển từ $n$ sang $n+1$, trong toàn bộ chuỗi bit của mã Gray, \textbf{chỉ có duy nhất một bit tại vị trí $k$ (chính là số lượng bit 1 tận cùng của $n$) bị lật đảo}, khoảng cách Hamming giữa $G(n)$ và $G(n+1)$ luôn luôn bằng 1:
\begin{equation}
\text{HammingDistance}\big(G(n), G(n+1)\big) \equiv 1 \quad \forall n \ge 0
\end{equation}

Tính chất này đã được kịch bản \texttt{verify\_v3/ch02\_gray\_code.py} kiểm chứng tự động cho toàn bộ $1.024$ bước nhảy liên tiếp, xác nhận không có bất kỳ ngoại lệ nào.

\subsubsection*{Lựa chọn kiến trúc và khuyến nghị thực hành trên Quartus}
Từ các phân tích định lượng trên, quy chuẩn thực hành thiết kế FSM cho hệ thống được xác lập như sau:
\begin{enumerate}[leftmargin=1.5em, itemsep=2pt]
    \item \textbf{Mẫu viết mã (Coding Template):} Áp dụng mẫu chuẩn 2 khối \texttt{always} (hoặc 3 khối \texttt{always}) rõ ràng: một khối tuần tự chốt thanh ghi trạng thái \texttt{state\_reg <= state\_next}, một khối tổ hợp thuần túy giải mã trạng thái tiếp theo \texttt{state\_next}, và các biểu thức gán liên tục gán ngõ ra. Đây là mẫu hình chuẩn mà bộ trích xuất FSM của Quartus (State Machine Extractor) dễ dàng nhận diện và phân tích chính xác.
    \item \textbf{Thiết lập tổng hợp:} Để tùy chọn \texttt{State Machine Processing} ở chế độ mặc định \textbf{Auto}. Trình biên dịch Quartus sẽ tự động ánh xạ FSM sang cấu hình One-Hot nhằm tối ưu hóa $F_{\max}$.
    \item \textbf{Kiểm tra báo cáo FSM:} Sau khi biên dịch, người thiết kế mở báo cáo \emph{Analysis \& Synthesis $\to$ State Machines} trong Quartus để đọc bảng mã hóa thực tế mà công cụ đã áp dụng \textbf{(phải đo)}.
    \item \textbf{Tần số cực đại $F_{\max}$ thực tế:} Các giá trị $F_{\max}$ của FSM mã hóa Binary so với One-Hot là thông số \textbf{(phải đo)} thông qua công cụ Timing Analyzer (Slow 1100mV 85C Model).
\end{enumerate}

\textit{Tóm lại: Trên kiến trúc Cyclone V với 166.036 thanh ghi dồi dào, mã hóa One-Hot giải phóng độ trễ giải mã tổ hợp; việc hiểu rõ bản chất toán học từ lập luận chuỗi bit 1 của mã Gray đến cấu trúc logic của One-Hot giúp ta tự tin để Quartus ở chế độ Auto và kiểm tra xác nhận qua báo cáo State Machine.}

\subsection*{2.4. Khối đệm trượt (Skid Buffer): Cắt đứt đường phản áp tới hạn}
\addcontentsline{toc}{subsection}{2.4. Khối đệm trượt (Skid Buffer): Cắt đứt đường phản áp tới hạn}

\subsubsection*{Vấn đề thực tế: Vòng lặp phản áp tổ hợp làm suy giảm Fmax}
Khi thiết kế các chuỗi đường ống (pipeline) nhiều tầng theo chuẩn Avalon-ST, các khối chức năng được mắc nối tiếp nhau: đầu ra của khối này (source) nối vào đầu vào của khối tiếp theo (sink). 

Để dòng dữ liệu chuyển động về phía trước (forward path), người thiết kế dễ dàng chèn các tầng thanh ghi flip-flop trên đường \texttt{data} và \texttt{valid}. Tuy nhiên, thách thức kỹ thuật nan giải lại nằm ở \textbf{đường tín hiệu phản áp đi ngược chiều (backward path)}: tín hiệu \texttt{ready}.

Hãy xét một tầng đệm đường ống đơn giản thông thường:
\begin{lstlisting}[language=Verilog, caption={Tầng chốt đường ống đơn giản với đường ready tổ hợp thuần}]
always @(posedge clk) begin
    if (m_ready) begin
        m_data  <= s_data;
        m_valid <= s_valid;
    end
end
assign s_ready = m_ready | ~m_valid;
\end{lstlisting}

Nhìn vào dòng lệnh gán liên tục cuối cùng: \texttt{s\_ready = m\_ready | \~{}m\_valid}. Tín hiệu \texttt{s\_ready} ở đầu vào phụ thuộc trực tiếp vào tín hiệu \texttt{m\_ready} ở đầu ra thông qua một cổng logic tổ hợp OR. 

Nếu ta ghép 5 khối xử lý như vậy nối tiếp nhau (ví dụ: Unpacker $\to$ DC-Blocker $\to$ FIR $\to$ Window RAM $\to$ MAC), tín hiệu \texttt{ready} từ khối MAC cuối cùng sẽ phải chạy xuyên qua 5 cổng tổ hợp OR ngược dòng từ hạ nguồn về tận bộ đệm FIFO ở thượng nguồn trong \textbf{cùng một chu kỳ xung nhịp 20 ns}. Đường dây phản áp này đi qua nhiều vùng logic khác nhau trên chip, chịu độ trễ dây dẫn (routing delay) lớn và trở thành \textbf{đường tới hạn (critical path)} quyết định tần số hoạt động cực đại $F_{\max}$ của cả hệ thống.

\textbf{Vì sao không thể chốt \texttt{ready} bằng một thanh ghi đơn giản?}
Một trực giác tự nhiên là đặt thêm một thanh ghi flip-flop trên đường \texttt{ready}: \texttt{always @(posedge clk) s\_ready <= m\_ready;}. 

Nhưng giải pháp này sẽ phá vỡ giao thức bắt tay và làm mất dữ liệu:
\begin{itemize}[leftmargin=1.5em, itemsep=2pt]
    \item Giả sử ở chu kỳ $n$, phía hạ nguồn bất ngờ bị nghẽn và hạ \texttt{m\_ready} xuống 0.
    \item Do tín hiệu \texttt{ready} bị trễ 1 nhịp qua thanh ghi, ở chính chu kỳ $n$ phía thượng nguồn vẫn nhìn thấy \texttt{s\_ready} = 1.
    \item Phía thượng nguồn phát đi một mẫu dữ liệu mới $D_{\text{mới}}$ ở chu kỳ $n$.
    \item Sang chu kỳ $n+1$, mẫu $D_{\text{mới}}$ này bay tới cửa vào của tầng đệm, nhưng hạ nguồn vẫn đang nghẽn (\texttt{m\_ready} = 0) và thanh ghi chính của tầng đệm đang bận giữ mẫu cũ $D_{\text{cũ}}$. 
    \item Vì không có chỗ chứa, mẫu $D_{\text{mới}}$ sẽ \textbf{ghi đè lên mẫu cũ hoặc bị rơi mất}.
\end{itemize}

Hiện tượng mẫu dữ liệu tiếp tục bay tới thêm 1 nhịp sau khi tín hiệu dừng đã được phát đi được gọi là hiện tượng \textbf{trượt quán tính (skid)}. Khối mạch giải quyết bài toán này mang tên \textbf{Khối đệm trượt (Skid Buffer)}.

\subsubsection*{Bản chất của Skid Buffer: Hai tầng chốt bảo toàn 100\% thông lượng}
Bản chất của Skid Buffer là cung cấp vừa đúng \textbf{hai vị trí lưu trữ dữ liệu} độc lập:
\begin{enumerate}[leftmargin=1.5em, itemsep=2pt]
    \item \textbf{Thanh ghi chính (Main Register - \texttt{main\_data}):} Giữ mẫu dữ liệu đang được trình diện ra phía hạ nguồn (\texttt{m\_data}).
    \item \textbf{Thanh ghi trượt (Skid Register - \texttt{skid\_data}):} Đóng vai trò vùng đệm dự phòng khẩn cấp, chuyên hấp thụ đúng 1 mẫu dữ liệu đang bay tới khi hạ nguồn đột ngột dừng tiếp nhận.
\end{enumerate}

Nguyên lý ngắt đường trễ tổ hợp của Skid Buffer thể hiện ở chỗ:
\begin{equation}
\texttt{s\_ready} = \overline{\texttt{skid\_valid}}
\end{equation}
Tín hiệu \texttt{s\_ready} báo cho phía thượng nguồn biết tầng đệm còn nhận được dữ liệu hay không \textbf{chỉ phụ thuộc vào cờ thanh ghi nội bộ \texttt{skid\_valid}}, hoàn toàn không có cổng logic tổ hợp nào nối từ \texttt{m\_ready}. Đường phản áp đi ngược đã bị cắt đứt bởi một thanh ghi flip-flop. Đồng thời, đường dữ liệu đi xuôi \texttt{m\_data} và \texttt{m\_valid} cũng được xuất trực tiếp từ các thanh ghi nội bộ. Cả hai chiều giao tiếp đều được cô lập bởi các phần tử nhớ tuần tự.

\begin{figure}[htbp]
\centering
% chuong_v3/fig_ch02_skid_buffer.tex
% Architecture of the two-stage registered Avalon-ST Skid Buffer
\begin{tikzpicture}[>=stealth, font=\small, node distance=1.5cm]

    % Styles
    \tikzstyle{block} = [draw=blue!80!black, fill=blue!5, thick, rounded corners=4pt, 
                         inner sep=8pt, align=center, font=\footnotesize\bfseries]
    \tikzstyle{regbox} = [draw=teal!80!black, fill=teal!8, thick, rounded corners=3pt, 
                          inner sep=6pt, align=center, font=\footnotesize\bfseries]
    \tikzstyle{ctrlbox} = [draw=purple!80!black, fill=purple!5, thick, rounded corners=3pt, 
                           inner sep=6pt, align=center, font=\footnotesize\bfseries]
    \tikzstyle{muxbox} = [draw=orange!80!black, fill=orange!10, thick, 
                          inner sep=4pt, align=center, font=\scriptsize\bfseries]
    \tikzstyle{pathline} = [->, thick, color=gray!80!black]
    \tikzstyle{dataline} = [->, very thick, color=teal!80!black]
    \tikzstyle{ctrlline} = [->, thick, dashed, color=purple!80!black]
    \tikzstyle{portlabel} = [font=\footnotesize\ttfamily]

    % Boundary box (Skid Buffer Module)
    \draw[draw=blue!60!black, fill=blue!2, thick, rounded corners=8pt, dashed] 
        (-3.2, -3.2) rectangle (6.8, 3.2);
    \node[anchor=north west, font=\normalsize\bfseries, color=blue!80!black] 
        at (-3.0, 3.0) {Khối đệm trượt (Skid Buffer) hai tầng chốt};

    % Upstream Interface (Left)
    \node[portlabel, anchor=east] (in_data)  at (-3.5,  1.5) {s\_data[15:0]};
    \node[portlabel, anchor=east] (in_valid) at (-3.5,  0.2) {s\_valid};
    \node[portlabel, anchor=east] (out_ready) at (-3.5, -1.8) {s\_ready};

    % Downstream Interface (Right)
    \node[portlabel, anchor=west] (out_data)  at (7.1,  1.5) {m\_data[15:0]};
    \node[portlabel, anchor=west] (out_valid) at (7.1,  0.2) {m\_valid};
    \node[portlabel, anchor=west] (in_ready)  at (7.1, -1.8) {m\_ready};

    % Internal Components:
    % 1. Skid Data Register (Top Left)
    \node[regbox, minimum width=2.4cm, minimum height=1.1cm] (skid_reg) at (-0.8, 1.8) 
        {Thanh ghi trượt\\(skid\_data)};

    % 2. MUX for Main Register input
    \node[muxbox, minimum width=0.8cm, minimum height=1.6cm] (mux_main) at (1.4, 1.5) 
        {MUX\\2:1};

    % 3. Main Data Register (Top Right)
    \node[regbox, minimum width=2.4cm, minimum height=1.1cm] (main_reg) at (4.2, 1.5) 
        {Thanh ghi chính\\(main\_data)};

    % 4. Skid Control Logic / Valid Register
    \node[ctrlbox, minimum width=2.6cm, minimum height=1.0cm] (skid_ctrl) at (-0.8, -0.6) 
        {Trạng thái trượt\\(skid\_valid FF)};

    % 5. Main Control Logic / Valid Register
    \node[ctrlbox, minimum width=2.6cm, minimum height=1.0cm] (main_ctrl) at (4.2, -0.6) 
        {Trạng thái chính\\(main\_valid FF)};

    % 6. Inverter for s_ready
    \node[draw=purple!80!black, fill=purple!15, circle, inner sep=2pt, font=\scriptsize\bfseries] (inv) at (-2.0, -1.8) {$\sim$};

    % Connections - Data path
    % s_data to Skid Register and MUX input 0
    \draw[dataline] (-3.2, 1.5) -- (-0.8, 1.5) |- (skid_reg.west);
    \draw[dataline] (-1.8, 1.5) |- (0.2, 1.2) -- (mux_main.west |- 0, 1.2);
    \node[font=\tiny, anchor=south] at (-2.2, 1.5) {dữ liệu vào};

    % Skid register to MUX input 1
    \draw[dataline] (skid_reg.east) -- (0.6, 1.8) |- (mux_main.west |- 0, 1.8);

    % MUX to Main Register
    \draw[dataline] (mux_main.east) -- (main_reg.west);

    % Main Register to m_data
    \draw[dataline] (main_reg.east) -- (6.8, 1.5);
    \node[font=\tiny, anchor=south] at (5.7, 1.5) {chốt ra};

    % Connections - Control path
    % s_valid in
    \draw[ctrlline] (-3.2, 0.2) -- (-1.8, 0.2) |- (skid_ctrl.west);
    \draw[ctrlline] (-1.8, 0.2) |- (2.2, 0.2) |- (main_ctrl.west);

    % m_ready in
    \draw[ctrlline] (6.8, -1.8) -- (4.2, -1.8) -- (main_ctrl.south);
    \draw[ctrlline] (4.2, -1.8) -| (0.2, -1.8) -- (skid_ctrl.south);

    % main_valid to output
    \draw[ctrlline] (main_ctrl.east) -- (6.8, -0.6) |- (6.8, 0.2);

    % skid_valid drives MUX select and s_ready inverter
    \draw[ctrlline] (skid_ctrl.north) -- (-0.8, 0.5) -| (mux_main.south);
    \node[font=\tiny, color=purple!80!black] at (0.8, 0.4) {sel=skid\_valid};

    \draw[ctrlline] (skid_ctrl.south) -- (-0.8, -1.8) -- (inv.east);
    \draw[ctrlline] (inv.west) -- (-3.2, -1.8);
    \node[font=\tiny, anchor=south, color=purple!80!black] at (-2.6, -1.7) {s\_ready};

    % Insight text at bottom
    \node[anchor=south, font=\scriptsize\itshape, color=gray!80!black] at (1.8, -3.0) 
        {* Tín hiệu s\_ready chỉ phụ thuộc vào FF nội bộ skid\_valid, triệt tiêu đường phản áp tổ hợp ngược.};

\end{tikzpicture}

\caption{Kiến trúc phần cứng của khối đệm trượt (Skid Buffer) hai tầng chốt: Đường phản áp \texttt{s\_ready} được điều khiển hoàn toàn bởi cờ nội bộ \texttt{skid\_valid}, ngắt đứt đường trễ tổ hợp giữa hai miền giao tiếp.}
\label{fig:ch02_skid_buffer}
\end{figure}

\textbf{Đọc hình.} Quan sát Hình~\ref{fig:ch02_skid_buffer}:
\begin{itemize}[leftmargin=1.5em, itemsep=2pt]
    \item Đường dữ liệu (vạch đậm màu xanh mòng két) đi vào từ \texttt{s\_data[15:0]}. Nó được rẽ làm hai nhánh: một nhánh đi thẳng vào bộ ghép kênh MUX 2:1 để nạp vào Thanh ghi chính, nhánh còn lại đi vào Thanh ghi trượt (\texttt{skid\_data}).
    \item Khi hạ nguồn sẵn sàng (\texttt{m\_ready} = 1) và thanh ghi trượt rỗng (\texttt{skid\_valid} = 0), dữ liệu đi thẳng qua MUX vào Thanh ghi chính mỗi chu kỳ xung nhịp, duy trì thông lượng liên tục $100\%$ ($1\text{ mẫu/nhịp}$).
    \item Khi hạ nguồn tạm dừng (\texttt{m\_ready} = 0) trong lúc thượng nguồn vẫn phát dữ liệu (\texttt{s\_valid} = 1), Thanh ghi chính bị khóa để bảo toàn dữ liệu cũ, và mẫu dữ liệu mới lập tức trượt vào lưu trú an toàn trong Thanh ghi trượt (\texttt{skid\_data}). Đồng thời, cờ \texttt{skid\_valid} bật lên 1, kích hoạt cổng đảo hạ \texttt{s\_ready} xuống 0 để chặn đứng thượng nguồn ở chu kỳ kế tiếp.
\end{itemize}

\subsubsection*{Ví dụ tính tay 2 nhịp trượt (Hand Trace Walkthrough)}
Theo dõi vết hoạt động chi tiết của Skid Buffer qua 5 chu kỳ xung nhịp liên tiếp với dòng dữ liệu gồm 3 mẫu $D_1, D_2, D_3$:

\begin{table}[htbp]
\centering
\footnotesize
\setlength{\tabcolsep}{3.5pt}
\caption{Bảng tính tay từng nhịp hoạt động của khối Skid Buffer khi xuất hiện phản áp}
\label{tab:ch02_skid_trace}
\begin{tabular}{cccccccp{4.8cm}}
\toprule
CK & \texttt{s\_dat} & \texttt{s\_vld} & \texttt{m\_rdy} & \texttt{main\_dat} & \texttt{skid\_dat} & \texttt{sk\_vld} & Hành động của mạch \\
\midrule
0 & $D_1$ & 1 & 1 & Trống & Trống & 0 & Mạch rỗng. $D_1$ chuẩn bị nạp vào Main. \\
1 & $D_2$ & 1 & \textbf{0} & $D_1$ & Trống & 0 & Hạ nguồn tắc (\texttt{m\_rdy}=0). $D_2$ bay tới! \\
2 & $D_3$ & 1 & \textbf{0} & $D_1$ & \textbf{$D_2$} & \textbf{1} & $D_2$ trượt vào Skid. \texttt{s\_ready}=0 chặn nguồn! \\
3 & $D_3$ & 1 & \textbf{1} & \textbf{$D_2$} & Trống & \textbf{0} & Hạ nguồn thông: $D_1$ xuất, $D_2$ sang Main! \\
4 & -- & 0 & 1 & $D_3$ & Trống & 0 & $D_2$ xuất, $D_3$ vào Main. Hoàn tất. \\
\bottomrule
\end{tabular}
\end{table}

\begin{itemize}[leftmargin=1.5em, itemsep=2pt]
    \item \textbf{Ở chu kỳ 1:} Thanh ghi chính đang giữ mẫu $D_1$ và chuẩn bị xuất ra. Bất ngờ hạ nguồn báo bận (\texttt{m\_ready} = 0). Cùng lúc đó, thượng nguồn đẩy mẫu $D_2$ lên bus với \texttt{s\_valid} = 1 (vì ở CK 1 \texttt{s\_ready} đang bằng 1).
    \item \textbf{Ở chu kỳ 2:} Đến cạnh xung nhịp của CK 2, do $D_1$ chưa được hạ nguồn tiếp nhận nên nó tiếp tục nằm lại trên Thanh ghi chính. Mẫu $D_2$ không thể vào Thanh ghi chính nên lập tức được chốt vào Thanh ghi trượt (\texttt{skid\_data} = $D_2$, \texttt{skid\_valid} = 1). Khi \texttt{skid\_valid} lên 1, \texttt{s\_ready} lập tức bị kéo xuống 0. Thượng nguồn dù có muốn gửi tiếp $D_3$ cũng bị chặn lại. Không có mẫu nào bị ghi đè.
    \item \textbf{Ở chu kỳ 3:} Hạ nguồn thông tắc (\texttt{m\_ready} = 1). Mẫu $D_1$ được chuyển giao thành công. Ngay lập tức, mạch điều khiển chuyển mẫu $D_2$ từ Thanh ghi trượt sang nạp đè vào Thanh ghi chính, và xóa cờ \texttt{skid\_valid} về 0. Cờ \texttt{s\_ready} tự động bật lên 1 trở lại, cho phép mẫu $D_3$ được nạp vào ở chu kỳ tiếp theo.
\end{itemize}

\subsubsection*{So sánh các phương án và lựa chọn thiết kế}
\begin{itemize}[leftmargin=1.5em, itemsep=2pt]
    \item \textbf{Phương án 1: Nối thẳng tổ hợp (Direct Combinational Ready):} Không tốn thêm thanh ghi, nhưng đường phản áp tổ hợp xuyên suốt nhiều khối làm giảm tần số hoạt động cực đại $F_{\max}$ \textbf{(phải đo)}. Chỉ phù hợp với các hệ thống có ít tầng xử lý và khoảng cách vật lý ngắn.
    \item \textbf{Phương án 2: Chèn Skid Buffer hai tầng chốt:} Trong mã nguồn\\ \texttt{rtl/v3/ch02\_skid\_buffer.v}, cấu trúc gồm 2 thanh ghi dữ liệu 16-bit (\texttt{main\_data}, \texttt{skid\_data}) và 2 cờ 1-bit (\texttt{main\_valid}, \texttt{skid\_valid}), tổng cộng gồm $16 \times 2 + 1 \times 2 = \textbf{34 flip-flop}$ khai báo ở mức RTL (số lượng thanh ghi trong mã nguồn; tài nguyên thực tế sau đóng gói ALM do Quartus Fitter báo cáo và phải đo). Thiết kế không dùng M10K và duy trì 100\% thông lượng (1 từ/chu kỳ).
\end{itemize}

\textbf{Lựa chọn của tài liệu:} Ta chọn \textbf{Phương án 2 (Skid Buffer)} tại các điểm giao tiếp phân chia miền chức năng quan trọng (như giữa bộ giải gói Unpacker và bộ nhớ RAM khung) để đảm bảo hệ thống đạt $F_{\max} \ge 50\text{ MHz}$ mà không phụ thuộc vào kết quả định tuyến ngẫu nhiên của phần mềm Quartus Fitter.

\subsubsection*{Đường (A): Khối IP Intel Avalon-ST Pipeline Stage}
Trong bộ công cụ thiết kế hệ thống Platform Designer (Qsys), tập đoàn Intel cung cấp sẵn một khối IP có chức năng tương tự mang tên \textbf{Avalon-ST Pipeline Stage}. 

Các tham số cấu hình chính trong Parameter Editor của IP gồm:
\begin{itemize}[leftmargin=1.5em, itemsep=2pt]
    \item \texttt{Data Width}: Bề rộng dữ liệu (chọn 16 bit cho mẫu âm thanh PCM).
    \item \texttt{Ready Latency}: Độ trễ sẵn sàng (chọn 0 chu kỳ).
    \item \texttt{Pipeline Type}: Kiểu đường ống. IP cho phép chọn một trong ba chế độ:
    \begin{enumerate}
        \item \emph{Forward}: Chỉ chốt thanh ghi trên đường \texttt{data} và \texttt{valid} (đường \texttt{ready} vẫn là tổ hợp).
        \item \emph{Backward}: Chỉ chốt thanh ghi trên đường \texttt{ready} (đòi hỏi xử lý trượt 1 nhịp).
        \item \emph{Both}: Chốt thanh ghi trên cả hai chiều. Khi chọn chế độ này, mã nguồn bên trong của IP sẽ tự động tổng hợp ra kiến trúc mạch đệm trượt hai vị trí lưu trữ tương đương với Skid Buffer.
    \end{enumerate}
\end{itemize}

Hạn chế của việc sử dụng IP hộp đen là cấu trúc sinh ra bởi Platform Designer thường kèm theo nhiều tín hiệu thừa (như \texttt{channel}, \texttt{error}, \texttt{empty}), gây cồng kềnh khi mô phỏng độc lập. Do đó, trong các thiết kế lõi DSP chuyên sâu, việc tự viết một mô-đun Skid Buffer chuẩn bằng Verilog-2001 (Đường B) luôn là lựa chọn ưu tiên hàng đầu của các kỹ sư FPGA.

\subsubsection*{Đường (B): Mã nguồn HDL tự viết chuẩn Verilog-2001}
Mã nguồn đầy đủ của khối đệm trượt được hiện thực hóa trong tệp \texttt{rtl/v3/ch02\_skid\_buffer.v}:

\begin{lstlisting}[language=Verilog, firstnumber=35, caption={Mô-đun Skid Buffer hai tầng chốt (rtl/v3/ch02\_skid\_buffer.v)}]
    // Main register: holds data currently presented to downstream
    reg [DATA_WIDTH-1:0] main_data;
    reg                  main_valid;

    // Skid register: holds in-flight data captured when downstream stalls
    reg [DATA_WIDTH-1:0] skid_data;
    reg                  skid_valid;

    // Upstream ready is asserted whenever the skid register is empty.
    // It is driven by a local flip-flop (skid_valid), with zero combinational
    // dependence on m_ready.
    assign s_ready = ~skid_valid;

    // Downstream outputs driven directly from main register
    assign m_valid = main_valid;
    assign m_data  = main_data;

    // Internal handshake conditions
    wire s_handshake = s_valid & s_ready;
    wire m_handshake = main_valid & m_ready;

    always @(posedge clk or negedge rst_n) begin
        if (!rst_n) begin
            main_data  <= {DATA_WIDTH{1'b0}};
            main_valid <= 1'b0;
            skid_data  <= {DATA_WIDTH{1'b0}};
            skid_valid <= 1'b0;
        end else begin
            if (m_handshake) begin
                // Downstream consumed current main_data
                if (skid_valid) begin
                    // Reload main from skid register; skid becomes empty
                    main_data  <= skid_data;
                    main_valid <= 1'b1;
                    skid_valid <= 1'b0;
                end else if (s_handshake) begin
                    // Direct feed from upstream to main register (full throughput streaming)
                    main_data  <= s_data;
                    main_valid <= 1'b1;
                end else begin
                    // Upstream had no new data; main register empties
                    main_valid <= 1'b0;
                end
            end else begin
                // Downstream stalled (m_ready == 0 or main_valid == 0)
                if (s_handshake) begin
                    if (main_valid) begin
                        // Main register is holding stalled data, so incoming beat skids into skid register
                        skid_data  <= s_data;
                        skid_valid <= 1'b1;
                    end else begin
                        // Main was empty; incoming beat fills main register
                        main_data  <= s_data;
                        main_valid <= 1'b1;
                    end
                end
            end
        end
    end
\end{lstlisting}

\textbf{Giải thích mã nguồn dòng-từng-dòng:}
\begin{itemize}[leftmargin=1.5em, itemsep=2pt]
    \item Dòng 36--41: Khai báo 4 thanh ghi nội bộ đại diện cho 2 vị trí lưu trữ: \texttt{main\_data} (16-bit) và \texttt{main\_valid} (1-bit); \texttt{skid\_data} (16-bit) và \texttt{skid\_valid} (1-bit).
    \item Dòng 46: \texttt{assign s\_ready = \~{}skid\_valid;}. Phía thượng nguồn được phép gửi dữ liệu khi và chỉ khi thanh ghi trượt đang rỗng. Biểu thức này không chứa biến \texttt{m\_ready}, chứng minh mạch đã triệt tiêu hoàn toàn đường trễ phản áp tổ hợp ngược.
    \item Dòng 49--50: Xuất dữ liệu và cờ hợp lệ ra hạ nguồn từ thanh ghi chính (\texttt{main\_data}, \texttt{main\_valid}).
    \item Dòng 53--54: Định nghĩa hai điều kiện bắt tay nội bộ \texttt{s\_handshake} và \texttt{m\_handshake}.
    \item Dòng 56--62: Nhánh reset bất đồng bộ tích cực mức thấp, đưa cả 4 thanh ghi về trạng thái rỗng ban đầu.
    \item Dòng 63--78: Trường hợp hạ nguồn hoàn tất một giao dịch (\texttt{m\_handshake} = 1). Nếu trong thanh ghi trượt đang có sẵn mẫu dữ liệu lưu cữu (\texttt{skid\_valid} = 1), mẫu đó lập tức được nạp sang \texttt{main\_data} và \texttt{skid\_valid} được xóa về 0. Nếu không có dữ liệu trượt nhưng thượng nguồn đang đưa mẫu mới vào (\texttt{s\_handshake} = 1), mẫu mới được nạp thẳng vào \texttt{main\_data}. Nếu cả hai nguồn đều không có dữ liệu, thanh ghi chính rỗng (\texttt{main\_valid} $\le$ 0).
    \item Dòng 79--91: Trường hợp hạ nguồn tắc nghẽn (\texttt{m\_handshake} = 0). Nếu thượng nguồn gửi dữ liệu tới (\texttt{s\_handshake} = 1): khi thanh ghi chính đã đầy (\texttt{main\_valid} = 1), mẫu mới bắt buộc phải rẽ nhánh vào chốt trong thanh ghi trượt (\texttt{skid\_data} $\le$ \texttt{s\_data}) và cờ \texttt{skid\_valid} bật lên 1. Nếu thanh ghi chính đang rỗng, dữ liệu được nạp vào thanh ghi chính.
\end{itemize}

\textit{Tóm lại: Khối đệm trượt (Skid Buffer) giải quyết mâu thuẫn giữa độ trễ tổ hợp của đường phản áp và yêu cầu toàn vẹn dữ liệu; với đúng hai tầng thanh ghi trong silicon, nó cho phép chốt chặn cả hai chiều phát-nhận mà vẫn duy trì thông lượng 1 mẫu/chu kỳ không một nhịp trễ.}

\subsection*{2.5. Kiểm chứng bắt tay luồng và thực nghiệm mã hóa FSM}
\addcontentsline{toc}{subsection}{2.5. Kiểm chứng bắt tay luồng và thực nghiệm mã hóa FSM}

\subsubsection*{Chiến lược kiểm chứng ba tầng (Three-Tier Verification)}
Để đảm bảo toàn bộ mã nguồn phần cứng của Chương 2 hoạt động chính xác trên silicon thực tế, quy trình kiểm chứng được tổ chức theo phương pháp luận ba tầng:
\begin{enumerate}[leftmargin=1.5em, itemsep=2pt]
    \item \textbf{Tầng 1 (Mô hình song sinh phần mềm Python Twin):} Xây dựng các lớp đối tượng Python mô phỏng chính xác hành vi của mạch theo từng chu kỳ xung nhịp và từng cạnh tích cực (\texttt{posedge clk}). Mô hình này chạy các kịch bản kích thích phức tạp (nghẽn ngẫu nhiên, sốc tải, xung reset đột ngột) và ghi nhận toàn bộ giá trị bus kỳ vọng ra tệp văn bản.
    \item \textbf{Tầng 2 (Mô phỏng chu kỳ HDL bằng Icarus Verilog):} Các testbench Verilog-2001 đọc trực tiếp tệp kích thích từ Tầng 1, cấp vào cổng của module phần cứng (DUT -- Device Under Test), đọc kết quả đầu ra sau mỗi nhịp xung nhịp và đối soát bit-đối-bit, chu kỳ-đối-chu kỳ với tệp kỳ vọng. Nếu có bất kỳ sự sai lệch nào dù chỉ 1 bit, testbench phát cờ \texttt{[FAIL]}.
    \item \textbf{Tầng 3 (Biên dịch Quartus và đo đạc phần cứng):} Nạp thiết kế qua Quartus Fitter để kiểm tra tính khả tổng hợp, đọc báo cáo State Machine Report và đo đạc tần số $F_{\max}$ thực tế trên bo mạch DE10-Standard.
\end{enumerate}

\subsubsection*{Kiểm chứng khối đệm trượt Skid Buffer với phản áp ngẫu nhiên}
Mô hình song sinh của khối Skid Buffer được hiện thực trong tệp \texttt{verify\_v3/ch02\_skid\_twin.py}. Dưới đây là hàm cập nhật chu kỳ trọng tâm của mô hình:

\begin{lstlisting}[language=Python, firstnumber=37, caption={Hàm cập nhật chu kỳ của mô hình song sinh Skid Buffer (verify\_v3/ch02\_skid\_twin.py)}]
    def step(self, rst_n, s_valid, s_data, m_ready):
        """
        Simulate one clock edge of rtl/v3/ch02_skid_buffer.v.
        Returns (s_ready, m_valid, m_data, skid_valid) observed AFTER posedge clk.
        """
        if not rst_n:
            self.reset()
            # In reset: main_valid=0, skid_valid=0 -> s_ready=1
            return (1, 0, 0, 0)
            
        s_data = s_data & self.mask
        
        # State before clock edge determines handshakes
        s_ready_before = 1 if (self.skid_valid == 0) else 0
        m_valid_before = self.main_valid
        
        s_handshake = (s_valid == 1) and (s_ready_before == 1)
        m_handshake = (m_valid_before == 1) and (m_ready == 1)
        
        next_main_data = self.main_data
        next_main_valid = self.main_valid
        next_skid_data = self.skid_data
        next_skid_valid = self.skid_valid
        
        if m_handshake:
            if self.skid_valid:
                next_main_data = self.skid_data
                next_main_valid = 1
                next_skid_valid = 0
            elif s_handshake:
                next_main_data = s_data
                next_main_valid = 1
            else:
                next_main_valid = 0
        else:
            if s_handshake:
                if self.main_valid:
                    next_skid_data = s_data
                    next_skid_valid = 1
                else:
                    next_main_data = s_data
                    next_main_valid = 1
                    
        # Register update on clock edge
        self.main_data = next_main_data
        self.main_valid = next_main_valid
        self.skid_data = next_skid_data
        self.skid_valid = next_skid_valid
        
        # Output signals after clock edge
        s_ready_after = 1 if (self.skid_valid == 0) else 0
        m_valid_after = self.main_valid
        m_data_after = self.main_data
        skid_valid_after = self.skid_valid
        
        return (s_ready_after, m_valid_after, m_data_after, skid_valid_after)
\end{lstlisting}

\textbf{Giải thích mã nguồn Python Twin dòng-từng-dòng:}
\begin{itemize}[leftmargin=1.5em, itemsep=2pt]
    \item Dòng 37: Phương thức \texttt{step(rst\_n, s\_valid, s\_data, m\_ready)} mô phỏng biến đổi trạng thái tại một chu kỳ xung nhịp.
    \item Dòng 42--45: Xử lý reset bất đồng bộ: khi \texttt{rst\_n == 0}, gọi \texttt{self.reset()} và trả về trạng thái rỗng (\texttt{s\_ready=1, m\_valid=0, m\_data=0, skid\_valid=0}).
    \item Dòng 50--51: Đọc trạng thái cờ trước cạnh xung nhịp: \texttt{s\_ready\_before} bằng 1 khi \texttt{skid\_valid == 0}; \texttt{m\_valid\_before} là giá trị \texttt{main\_valid} hiện thời.
    \item Dòng 53--54: Đánh giá hai điều kiện bắt tay \texttt{s\_handshake} và \texttt{m\_handshake}.
    \item Dòng 56--59: Thiết lập các biến tạm lưu giá trị kế tiếp của 4 thanh ghi.
    \item Dòng 61--70: Xử lý khi hạ nguồn nhận mẫu (\texttt{m\_handshake} đúng): ưu tiên nạp mẫu từ thanh ghi trượt nếu \texttt{skid\_valid == 1}; nếu không có dữ liệu trượt thì nạp từ thượng nguồn; nếu cả hai đều không có dữ liệu thì hạ \texttt{main\_valid} về 0.
    \item Dòng 71--79: Xử lý khi hạ nguồn bận (\texttt{m\_handshake} sai): nếu thượng nguồn gửi mẫu (\texttt{s\_handshake} đúng), khi \texttt{main\_valid} đang bận, mẫu mới trượt vào \texttt{skid\_data} và bật cờ \texttt{skid\_valid} lên 1; nếu thanh ghi chính rỗng, nạp vào thanh ghi chính.
    \item Dòng 81--84: Cập nhật các trường thanh ghi của đối tượng tại cạnh lên xung nhịp.
    \item Dòng 87--92: Trích xuất và trả về bộ 4 tín hiệu ngõ ra quan sát được sau cạnh xung nhịp để phục vụ đối soát với testbench Verilog.
\end{itemize}

Dưới đây là bảng đối chiếu ánh xạ trực tiếp giữa mô hình Python Twin và mã nguồn phần cứng HDL:

\begin{table}[htbp]
\centering
\footnotesize
\caption{Ánh xạ tương đương giữa Python Twin và mã nguồn HDL của Skid Buffer}
\label{tab:ch02_skid_mapping}
\begin{tabularx}{\textwidth}{ccX}
\toprule
Dòng Twin (\texttt{ch02\_skid\_twin.py}) & Dòng HDL (\texttt{ch02\_skid\_buffer.v}) & Ý nghĩa kỹ thuật tương đương \\
\midrule
Dòng 42--45 & Dòng 57--62 & Xóa rỗng các thanh ghi khi reset tích cực (\texttt{rst\_n} = 0) \\
Dòng 50 & Dòng 46 & Xác định cờ sẵn sàng thượng nguồn \texttt{s\_ready = \~{}skid\_valid} \\
Dòng 53--54 & Dòng 53--54 & Đánh giá điều kiện hoàn tất bắt tay \texttt{s\_handshake} và \texttt{m\_handshake} \\
Dòng 61--70 & Dòng 63--78 & Xử lý khi hạ nguồn nhận mẫu: giải phóng trượt hoặc chuyển tiếp \\
Dòng 71--79 & Dòng 79--91 & Xử lý khi hạ nguồn nghẽn: hấp thụ mẫu đang bay vào thanh ghi trượt \\
Dòng 81--84 & Dòng 56--93 & Chốt cập nhật trạng thái thanh ghi tại cạnh lên xung nhịp (\texttt{posedge clk}) \\
Dòng 87--92 & Dòng 46--50 & Xuất giá trị tín hiệu giao tiếp sau cạnh xung nhịp để đối soát \\
\bottomrule
\end{tabularx}
\end{table}

Kịch bản mô phỏng tạo ra chuỗi kiểm thử dài \textbf{600 chu kỳ xung nhịp} liên tục, đi qua 6 pha thử thách:
\begin{itemize}[leftmargin=1.5em, itemsep=2pt]
    \item \textbf{Pha 1 (CK 0--4):} Xung reset tích cực mức thấp, kiểm tra trạng thái xóa rỗng ban đầu của 4 thanh ghi.
    \item \textbf{Pha 2 (CK 5--39):} Truyền dòng chảy tốc độ cao không nghẽn (\texttt{s\_valid} = 1, \texttt{m\_ready} = 1 liên tục 35 nhịp), kiểm tra thông lượng 1 mẫu/chu kỳ.
    \item \textbf{Pha 3 (CK 40--99):} Phản áp nhịp đơn định kỳ (cứ 3 chu kỳ hạ \texttt{m\_ready} xuống 0 một lần), kiểm tra khả năng kích hoạt và giải phóng thanh ghi trượt \texttt{skid\_data}.
    \item \textbf{Pha 4 (CK 100--179):} Phản áp dài hạn kéo dài liên tục 15 chu kỳ (tại các khoảng CK 120--135 và CK 150--165), kiểm tra khả năng khóa thượng nguồn \texttt{s\_ready} = 0 mà không làm tràn bộ nhớ.
    \item \textbf{Pha 5 (CK 180--549):} Bắt tay giả ngẫu nhiên (Pseudo-random backpressure với xác suất phát $70\%$ và xác suất nhận $60\%$), mô phỏng môi trường hoạt động thực tế khi kết nối giữa các khối DSP.
    \item \textbf{Pha 6 (CK 550--599):} Xả cạn bộ đệm về trạng thái nghỉ.
\end{itemize}

Kết quả chạy kịch bản Python twin ghi nhận: trong suốt 600 chu kỳ, phía thượng nguồn đã phát \textbf{292 mẫu âm thanh}, và phía hạ nguồn đã tiếp nhận chính xác \textbf{292/292 mẫu} (đạt tỷ lệ 100\%), không mất một mẫu nào, không lặp một mẫu nào, thứ tự thời gian được bảo toàn.

Testbench Verilog tương ứng \texttt{tb/v3/tb\_ch02\_skid\_buffer.v} thực hiện đối soát 600 chu kỳ này trên phần cứng RTL:
\begin{center}
\todo{box}
\end{center}

\subsubsection*{Kiểm chứng tương đương logic giữa FSM nhị phân và FSM One-Hot}
Để chứng minh hai phương án mã hóa FSM (Minimal Bits 3-bit trong \texttt{rtl/v3/ch02\_fsm\_binary.v} và One-Hot 5-bit trong \texttt{rtl/v3/ch02\_fsm\_onehot.v}) tương đương nhau về mặt hành vi chức năng, một testbench kiểm thử tương đương kép (dual-DUT testbench) được xây dựng trong tệp\\ \texttt{tb/v3/tb\_ch02\_fsm\_equiv.v}.

Mô hình song sinh Python Twin được cài đặt trong tệp \texttt{verify\_v3/ch02\_fsm\_twin.py}:

\begin{lstlisting}[language=Python, firstnumber=39, caption={Hàm cập nhật chu kỳ của mô hình song sinh FSM (verify\_v3/ch02\_fsm\_twin.py)}]
    def step(self, rst_n, start, load_done, calc_done, drain_done):
        if not rst_n:
            self.state = self.IDLE
            return (0, 0, 0, 0, 0, self.IDLE)
            
        curr_state = self.state
        
        # Next state logic
        next_state = curr_state
        if curr_state == self.IDLE:
            if start:
                next_state = self.LOAD
        elif curr_state == self.LOAD:
            if load_done:
                next_state = self.CALC
        elif curr_state == self.CALC:
            if calc_done:
                next_state = self.DRAIN
        elif curr_state == self.DRAIN:
            if drain_done:
                next_state = self.DONE
        elif curr_state == self.DONE:
            next_state = self.IDLE
            
        # Register updates on posedge clk
        self.state = next_state
        
        # Outputs observed after posedge clk (Moore outputs of current state)
        busy = 1 if self.state != self.IDLE else 0
        load_en = 1 if self.state == self.LOAD else 0
        calc_en = 1 if self.state == self.CALC else 0
        drain_en = 1 if self.state == self.DRAIN else 0
        done_pulse = 1 if self.state == self.DONE else 0
        state_num = self.state
        
        return (busy, load_en, calc_en, drain_en, done_pulse, state_num)
\end{lstlisting}

\textbf{Giải thích mã nguồn Python Twin FSM dòng-từng-dòng:}
\begin{itemize}[leftmargin=1.5em, itemsep=2pt]
    \item Dòng 39: Hàm \texttt{step(rst\_n, start, ...)} nhận các cờ kích hoạt rẽ nhánh và xung nhịp.
    \item Dòng 40--42: Nhánh reset bất đồng bộ đưa trạng thái về \texttt{IDLE} (0) và trả về toàn bộ tín hiệu đầu ra ở mức 0.
    \item Dòng 47--62: Logic giải mã chuyển trạng thái: từ \texttt{IDLE} sang \texttt{LOAD} khi có \texttt{start}; từ \texttt{LOAD} sang \texttt{CALC} khi có \texttt{load\_done}; từ \texttt{CALC} sang \texttt{DRAIN} khi có \texttt{calc\_done}; từ \texttt{DRAIN} sang \texttt{DONE} khi có \texttt{drain\_done}; từ \texttt{DONE} quay về \texttt{IDLE}.
    \item Dòng 64: Cập nhật biến trạng thái nội bộ \texttt{self.state = next\_state} tại cạnh lên xung nhịp.
    \item Dòng 67--72: Giải mã các tín hiệu ngõ ra Moore từ trạng thái \texttt{self.state}.
    \item Dòng 74: Trả về bộ 6 tín hiệu ngõ ra để đối soát với DUT phần cứng.
\end{itemize}

Dưới đây là bảng đối chiếu ánh xạ tương đương giữa Python Twin và hai module phần cứng:

\begin{table}[htbp]
\centering
\footnotesize
\caption{Ánh xạ tương đương giữa Python Twin và mã nguồn HDL FSM (Binary và One-Hot)}
\label{tab:ch02_fsm_mapping}
\begin{tabularx}{\textwidth}{ccX}
\toprule
Dòng Twin (\texttt{ch02\_fsm\_twin.py}) & Dòng HDL (\texttt{ch02\_fsm\_*}) & Ý nghĩa kỹ thuật tương đương \\
\midrule
Dòng 40--42 & Dòng 50--52 (Bin) / 51--53 (One) & Xóa thanh ghi trạng thái về \texttt{S\_IDLE} khi có tín hiệu reset \\
Dòng 47--62 & Dòng 58--80 (Bin) / 59--66 (One) & Logic tổ hợp giải mã trạng thái kế tiếp từ các cờ kích hoạt \\
Dòng 64 & Dòng 49--55 (Bin) / 50--56 (One) & Chốt chuyển trạng thái \texttt{state\_reg <= state\_next} tại \texttt{posedge clk} \\
Dòng 67--72 & Dòng 83--88 (Bin) / 69--73 (One) & Sinh các tín hiệu điều khiển Moore (\texttt{busy}, cờ \texttt{\_en}, \texttt{done\_pulse}) \\
Dòng 72 & Dòng 88 (Bin) / 76--85 (One) & Xuất chỉ số trạng thái số 3-bit \texttt{state\_num} để đối soát \\
\bottomrule
\end{tabularx}
\end{table}

Kịch bản kích thích \texttt{verify\_v3/ch02\_fsm\_twin.py} sinh ra \textbf{500 chu kỳ xung nhịp} với các tín hiệu điều khiển rẽ nhánh ngẫu nhiên (\texttt{start}, \texttt{load\_done}, \texttt{calc\_done}, \texttt{drain\_done}) và các xung reset bất thường ở chu kỳ 50 và 150.

Tại mỗi chu kỳ xung nhịp, testbench thực hiện hai bước đối soát đồng thời:
\begin{enumerate}[leftmargin=1.5em, itemsep=2pt]
    \item So sánh đối đầu trực tiếp giữa đầu ra của module Binary và module One-Hot: kiểm tra tính đồng nhất của toàn bộ 6 tín hiệu ngõ ra (\texttt{busy}, \texttt{load\_en}, \texttt{calc\_en}, \texttt{drain\_en}, \texttt{done\_pulse}, và mã trạng thái số \texttt{state\_num}).
    \item So sánh đầu ra của cả hai module với giá trị chuẩn từ mô hình song sinh Python\\ \texttt{verify\_v3/tb\_ch02\_fsm\_expected.txt}.
\end{enumerate}

Kết quả mô phỏng ghi nhận:
\begin{center}
\todo{box}
\end{center}

\subsubsection*{Bảng so sánh tổng hợp: Khối IP Intel so với HDL tự viết}
Dưới đây là bảng tổng kết định lượng đối chiếu giữa giải pháp sử dụng khối IP có sẵn của Intel và giải pháp tự thiết kế bằng HDL cho các thành phần của Chương 2:

\begin{table}[htbp]
\centering
\footnotesize
\setlength{\tabcolsep}{4.5pt}
\caption{So sánh toàn diện giữa khối IP Intel và mã nguồn HDL tự viết cho Chương 2}
\label{tab:ch02_ip_vs_hdl}
\begin{tabular}{lcc}
\toprule
Tiêu chí kỹ thuật & Khối IP Intel (Platform Designer) & HDL tự viết (Tài liệu này) \\
\midrule
Tên thành phần tương ứng & \texttt{altera\_avalon\_st\_pipeline\_stage} & \texttt{ch02\_skid\_buffer} \\
Tài nguyên khối nhớ M10K & 0 khối & 0 khối \\
Tài nguyên bộ nhân DSP & 0 khối & 0 khối \\
Tài nguyên thanh ghi (FF) & Phụ thuộc Qsys \textbf{(phải đo)} & \textbf{34 FF} (RTL: 32 data + 2 valid) \textbf{(phải đo)} \\
Tài nguyên khối logic ALM & Phụ thuộc tối ưu hóa Qsys \textbf{(phải đo)} & \textbf{Phải đo trong Quartus} \\
Thông lượng cực đại & 1 từ/chu kỳ (khi cấu hình \emph{Both}) & \textbf{1 từ/chu kỳ (100\%)} \\
Độ trễ tối thiểu (Latency) & 1 chu kỳ xung nhịp & 1 chu kỳ xung nhịp \\
Cắt đứt trễ tổ hợp \texttt{ready}? & Có (khi chọn chế độ \emph{Both}) & \textbf{Có (cắt đứt hoàn toàn)} \\
Tính độc lập và trong suốt & Phụ thuộc thư viện lõi Qsys & \textbf{Verilog-2001 chuẩn} \\
Khả năng kiểm thử chu kỳ & Khó kiểm thử độc lập ngoài Quartus & \textbf{Kiểm thử tức thì với iverilog} \\
Mã hóa FSM điều khiển & Thiết lập \texttt{Auto} trong Quartus & \textbf{Mã hóa tường minh One-Hot/Binary} \\
Báo cáo FSM trong Quartus & Xem State Machine Report \textbf{(phải đo)} & Xem State Machine Report \textbf{(phải đo)} \\
\bottomrule
\end{tabular}
\end{table}

\textbf{Đọc bảng.} Bảng~\ref{tab:ch02_ip_vs_hdl} chỉ rõ rằng: về mặt tài nguyên và hiệu năng lý thuyết, khối Skid Buffer tự viết bằng HDL tương đương với khối IP Pipeline Stage của Intel ở cấu hình chốt kép. Tuy nhiên, mã nguồn HDL tự viết mang lại lợi thế về tính trong suốt (transparency), không bị trói buộc vào môi trường sinh mã tự động của Platform Designer, cho phép lập trình viên kiểm thử từng bit độc lập với trình mô phỏng nhẹ Icarus Verilog và dễ dàng tích hợp vào bất kỳ họ FPGA nào.

\subsubsection*{Bài tập cuối chương có lời giải chi tiết}

\paragraph{Bài tập 1: Đánh giá độ trễ và thông lượng khi ghép tầng Skid Buffer.}
Giả sử trên đường truyền dữ liệu âm thanh từ bộ giải gói Unpacker tới bộ nhớ RAM khung, ta chèn thêm 2 khối Skid Buffer liên tiếp để cô lập định tuyến. Hãy tính:
\begin{enumerate}[label=(\alph*)]
    \item Độ trễ ban đầu (initial latency) tính bằng số chu kỳ xung nhịp từ khi mẫu đầu tiên được đưa vào cổng \texttt{s\_data} của Skid Buffer 1 đến khi nó xuất hiện tại cổng \texttt{m\_data} của Skid Buffer 2 trong điều kiện hạ nguồn luôn sẵn sàng (\texttt{m\_ready} = 1).
    \item Thông lượng ổn định (steady-state throughput) của hệ thống tính bằng số mẫu/giây ở xung nhịp 50 MHz.
\end{enumerate}

\textbf{Lời giải mẫu:}
\begin{enumerate}[label=(\alph*)]
    \item Theo phân tích ở Mục 2.4, mỗi khối Skid Buffer khi hoạt động ở trạng thái bình thường (không nghẽn) chốt dữ liệu qua một tầng thanh ghi chính (\texttt{main\_data}), đóng góp đúng 1 chu kỳ xung nhịp vào độ trễ đường ống. Khi ghép $P = 2$ khối Skid Buffer nối tiếp nhau, mẫu đầu tiên đi qua Skid Buffer 1 mất 1 chu kỳ và đi qua Skid Buffer 2 mất thêm 1 chu kỳ. Do đó, độ trễ ban đầu là:
    \begin{equation}
    L = 2 \times 1 = \textbf{2 chu kỳ xung nhịp} \quad (= 2 \times 20\text{ ns} = 40\text{ ns})
    \end{equation}
    \item Sau khi đường ống 2 tầng đã được nạp đầy (từ chu kỳ thứ 2 trở đi), tại mỗi nhịp xung nhịp 50 MHz đều có đúng 1 mẫu âm thanh xuất hiện ở đầu ra. Do đó, thông lượng ổn định là:
    \begin{equation}
    \Phi = 1\text{ mẫu/chu kỳ} \times 50.000.000\text{ chu kỳ/s} = \textbf{50.000.000 mẫu/giây}
    \end{equation}
    Thông lượng này vượt xa tốc độ yêu cầu của dòng âm thanh 16 kHz ($16.000\text{ mẫu/giây}$), chứng minh việc chèn thêm các tầng Skid Buffer hoàn toàn không làm suy giảm tốc độ truyền tải thực tế của hệ thống.
\end{enumerate}

\paragraph{Bài tập 2: Tính toán đường tới hạn của chuỗi phản áp tổ hợp khi không dùng Skid Buffer.}
Xét một chuỗi xử lý gồm $K = 6$ khối phần cứng mắc nối tiếp nhau theo chuẩn Avalon-ST. Các khối sử dụng tầng đệm đơn giản với đường \texttt{ready} tổ hợp thuần túy (\texttt{s\_ready = m\_ready | \~{}m\_valid}). Giả sử trong mỗi khối, cổng logic tổ hợp OR đóng góp thời gian trễ $t_{\text{gate}} = 0{,}8\text{ ns}$, và dây dẫn định tuyến giữa hai khối liên tiếp đóng góp thời gian trễ $t_{\text{wire}} = 1{,}6\text{ ns}$. Thời gian thiết lập (setup time) của thanh ghi nguồn là $t_{\text{su}} = 0{,}3\text{ ns}$ và thời gian trễ từ xung nhịp tới ngõ ra của thanh ghi đích là $t_{\text{co}} = 0{,}5\text{ ns}$.
\begin{enumerate}[label=(\alph*)]
    \item Hãy lập công thức và tính tổng thời gian trễ trên đường phản áp \texttt{ready} từ khối thứ 6 ngược về khối thứ 1.
    \item Xác định tần số hoạt động cực đại $F_{\max}$ lý thuyết của đường phản áp này. Liệu hệ thống có đáp ứng được xung nhịp 50 MHz hay không?
\end{enumerate}

\textbf{Lời giải mẫu:}
\begin{enumerate}[label=(\alph*)]
    \item Tín hiệu phản áp \texttt{ready} xuất phát từ thanh ghi của khối thứ 6 (đóng góp trễ $t_{\text{co}}$), đi qua $K = 6$ cổng tổ hợp OR và $K-1 = 5$ đoạn dây dẫn liên khối, trước khi đi vào thanh ghi điều khiển của khối thứ 1 (đòi hỏi thời gian thiết lập $t_{\text{su}}$). Tổng thời gian trễ trên đường tới hạn này là:
    \begin{align}
    T_{\text{crit}} &= t_{\text{co}} + (K \times t_{\text{gate}}) + \big((K - 1) \times t_{\text{wire}}\big) + t_{\text{su}} \nonumber \\
                    &= 0{,}5\text{ ns} + (6 \times 0{,}8\text{ ns}) + (5 \times 1{,}6\text{ ns}) + 0{,}3\text{ ns} \nonumber \\
                    &= 0{,}5 + 4{,}8 + 8{,}0 + 0{,}3 = \textbf{13{,}6 ns}
    \end{align}
    \item Chu kỳ xung nhịp tối thiểu cho phép là $T_{\min} = T_{\text{crit}} = 13{,}6\text{ ns}$. Tần số hoạt động cực đại tương ứng là:
    \begin{equation}
    F_{\max} = \frac{1}{T_{\min}} = \frac{1}{13{,}6 \times 10^{-9}\text{ s}} \approx \textbf{73{,}53 MHz}
    \end{equation}
    Ở tần số 50 MHz (chu kỳ $T_{\text{clk}} = 20\text{ ns}$), mạch vẫn thỏa mãn do $13{,}6\text{ ns} < 20\text{ ns}$ với khoảng dự trữ định thời (slack) là $20 - 13{,}6 = 6{,}4\text{ ns}$. Tuy nhiên, nếu tần số xung nhịp hệ thống nâng lên 100 MHz ($T_{\text{clk}} = 10\text{ ns}$) hoặc khi phần mềm Quartus định tuyến dây dẫn dài hơn ($t_{\text{wire}} > 2{,}8\text{ ns}$ do chip bị phân mảnh tài nguyên), mạch sẽ vi phạm định thời thời gian thiết lập (setup timing violation). Đây chính là lý do vì sao việc ngắt chuỗi phản áp bằng Skid Buffer là giải pháp an toàn trong thiết kế chuẩn công nghiệp.
\end{enumerate}

\paragraph{Bài tập 3: Dẫn xuất mã Gray 4-bit và kiểm chứng chuỗi bit 1 tận cùng.}
Hãy lập bảng chuyển đổi mã Gray 4-bit cho 8 số nguyên đầu tiên $n = 0..7$ bằng công thức $G(n) = n \oplus (n \gg 1)$. Chỉ rõ vị trí bit bị lật đảo tại mỗi bước chuyển và đối chiếu với số lượng bit 1 tận cùng ($k$) của số $n$.

\textbf{Lời giải mẫu:}
Áp dụng công thức $G(n) = n \oplus (n \gg 1)$ theo từng bước:
\begin{itemize}[leftmargin=1.5em, itemsep=2pt]
    \item $n = 0 = (0000)_2$: Không có bit 1 tận cùng ($k = 0$). $G(0) = 0000 \oplus 0000 = (0000)_2$.
    \item $n = 1 = (0001)_2$: Có 1 bit 1 tận cùng ($k = 1$). $G(1) = 0001 \oplus 0000 = (0001)_2$. Lệch với $G(0)$ ở bit 0 ($k=0$ của bước chuyển $0 \to 1$).
    \item $n = 2 = (0010)_2$: Có 0 bit 1 tận cùng ($k = 0$). $G(2) = 0010 \oplus 0001 = (0011)_2$. Lệch với $G(1)$ ở bit 1 ($k=1$ của số 1).
    \item $n = 3 = (0011)_2$: Có 2 bit 1 tận cùng ($k = 2$). $G(3) = 0011 \oplus 0001 = (0010)_2$. Lệch với $G(2)$ ở bit 0 ($k=0$ của số 2).
    \item $n = 4 = (0100)_2$: Có 0 bit 1 tận cùng ($k = 0$). $G(4) = 0100 \oplus 0010 = (0110)_2$. Lệch với $G(3)$ ở bit 2 ($k=2$ của số 3).
    \item $n = 5 = (0101)_2$: Có 1 bit 1 tận cùng ($k = 1$). $G(5) = 0101 \oplus 0010 = (0111)_2$. Lệch với $G(4)$ ở bit 0 ($k=0$ của số 4).
    \item $n = 6 = (0110)_2$: Có 0 bit 1 tận cùng ($k = 0$). $G(6) = 0110 \oplus 0011 = (0101)_2$. Lệch với $G(5)$ ở bit 1 ($k=1$ của số 5).
    \item $n = 7 = (0111)_2$: Có 3 bit 1 tận cùng ($k = 3$). $G(7) = 0111 \oplus 0011 = (0100)_2$. Lệch với $G(6)$ ở bit 0 ($k=0$ của số 6).
\end{itemize}
Tại mọi bước chuyển $n \to n+1$, vị trí bit bị đổi trạng thái luôn trùng khớp chính xác với số lượng bit 1 tận cùng của số $n$, khoảng cách Hamming luôn bằng 1, hoàn toàn đúng với định lý đã chứng minh ở Mục 2.3.

\paragraph{Bài tập 4: Đánh giá chi phí phần cứng của FSM điều khiển tự tương quan.}
Máy trạng thái điều khiển vòng lặp tự tương quan MAC trong Chương 7 có 8 trạng thái hoạt động: \texttt{RESET}, \texttt{WAIT\_HOP}, \texttt{INIT\_R0}, \texttt{CALC\_R0}, \texttt{NEXT\_LAG}, \texttt{CALC\_LAG}, \texttt{STORE\_PEAK}, và \texttt{EMIT\_FRAME}. Hãy so sánh:
\begin{enumerate}[label=(\alph*)]
    \item Số lượng thanh ghi flip-flop cần thiết giữa phương án Minimal Bits và One-Hot.
    \item Biểu thức logic kích hoạt tín hiệu nạp tích lũy \texttt{accum\_sload} (chỉ tích cực tại hai trạng thái \texttt{INIT\_R0} và \texttt{NEXT\_LAG}) trong cả hai phương án mã hóa. Nhận xét về độ trễ giải mã.
\end{enumerate}

\textbf{Lời giải mẫu:}
\begin{enumerate}[label=(\alph*)]
    \item Với $S = 8$ trạng thái:
    \begin{itemize}
        \item Phương án Minimal Bits: Số flip-flop là $K_{\text{bin}} = \lceil \log_2 8 \rceil = \textbf{3 thanh ghi FF}$.
        \item Phương án One-Hot: Số flip-flop là $K_{\text{one}} = 8 = \textbf{8 thanh ghi FF}$.
    \end{itemize}
    \item Biểu thức logic kích hoạt tín hiệu \texttt{accum\_sload}:
    \begin{itemize}
        \item Trong phương án Minimal Bits: Giả sử \texttt{INIT\_R0} được gán mã nhị phân \texttt{010} ($2$) và \texttt{NEXT\_LAG} được gán mã \texttt{100} ($4$). Biểu thức giải mã tổ hợp là:
        \begin{equation}
        \texttt{accum\_sload}_{\text{bin}} = (\overline{s_2} \cdot s_1 \cdot \overline{s_0}) + (s_2 \cdot \overline{s_1} \cdot \overline{s_0})
        \end{equation}
        Tín hiệu này cần một bảng tra LUT 3 đầu vào, đóng góp trễ truyền lan khoảng $0{,}3\text{ ns}$ đến $0{,}5\text{ ns}$ trước khi tới cổng điều khiển của khối DSP.
        \item Trong phương án One-Hot: Đặt $s[2]$ là bit của \texttt{INIT\_R0} và $s[4]$ là bit của \texttt{NEXT\_LAG}. Biểu thức logic là:
        \begin{equation}
        \texttt{accum\_sload}_{\text{one}} = s[2] + s[4]
        \end{equation}
        Biểu thức này chỉ là một cổng OR 2 đầu vào đơn giản (hoặc được tối ưu trực tiếp vào chân kích hoạt của DSP), cho độ trễ rất nhỏ, giúp đảm bảo tín hiệu điều khiển \texttt{accum\_sload} đến kịp thời cùng lúc với toán hạng nhân ở cạnh xung nhịp 50 MHz.
    \end{itemize}
\end{enumerate}

\textit{Tóm lại: Kiểm chứng chu kỳ với mô hình song sinh phần mềm và kịch bản kích thích ngẫu nhiên trên hơn 500 vector xác nhận rằng khối đệm trượt bảo toàn 100\% dòng dữ liệu không mất mát, trong khi hai lối kiến trúc FSM nhị phân và one-hot cho kết quả logic tương đương theo từng nhịp xung nhịp.}
